\documentclass[12pt,openany]{ctexbook}

% 上海财经大学硕士学位论文版式
% 版式目标：A4，黑白色正文，左侧装订，三线表，表题在上，图题在下
% 生效源文件：chapters_new.tex（第1-5章）、chapters_6_9.tex（第6-7章与附录）、backmatter.tex。
% chapter7_new.tex、chapter89_new.tex、appendix_new.tex 不参与编译，仅作历史存档。
\usepackage[a4paper,top=3.8cm,bottom=3.8cm,left=3.2cm,right=3.2cm,headheight=13pt,headsep=0.35cm,footskip=0.8cm]{geometry}
\usepackage{setspace}
\usepackage{amsmath,amssymb,amsthm}
\usepackage{booktabs}
\usepackage{graphicx}
\usepackage{enumitem}
\usepackage{array}
\usepackage{caption}
\usepackage{tocloft}
\renewcommand{\cftchapfont}{\mdseries\songti\zihao{4}}
\renewcommand{\cftchappagefont}{\mdseries\songti\zihao{4}}
\renewcommand{\cftsecfont}{\mdseries\songti\zihao{-4}}
\renewcommand{\cftsecpagefont}{\mdseries\songti\zihao{-4}}
\renewcommand{\cftsubsecfont}{\mdseries\songti\zihao{5}}
\renewcommand{\cftsubsecpagefont}{\mdseries\songti\zihao{5}}
\setlength{\cftchapindent}{2em}
\setlength{\cftsecindent}{4em}
\setlength{\cftsecnumwidth}{4.2em}
\renewcommand{\cftsecaftersnum}{\hspace{1em}}
\setlength{\cftsubsecindent}{7em}
\setlength{\cftsubsecnumwidth}{2.5em}
\renewcommand{\cftsubsecaftersnum}{\hspace{0.5em}}
\usepackage{fancyhdr}
\usepackage[hidelinks]{hyperref}
\usepackage[round,authoryear]{natbib}
\bibpunct{(}{)}{;}{a}{,}{,}
\renewcommand{\bibname}{参考文献}
\usepackage{fontspec}
\usepackage{seqsplit}
\usepackage{tikz}
\usepackage{algorithm}
\usepackage{float}
\usepackage{placeins}
\usepackage{algpseudocode}
\usepackage{listings}
\lstset{
  basicstyle=\ttfamily\footnotesize,
  breaklines=true,
  columns=fullflexible,
  frame=single,
  numbers=left,
  numberstyle=\tiny,
  showstringspaces=false,
  keepspaces=true,
  tabsize=2
}
\usetikzlibrary{arrows.meta,positioning,fit,backgrounds}

\newfontfamily\arial{Arial}
\setmainfont{Times New Roman}
\setCJKmainfont[AutoFakeBold=2,BoldFont=SimHei,ItalicFont=SimSun]{SimSun}
\setCJKsansfont[AutoFakeBold=2,BoldFont=SimHei]{SimHei}
\setCJKmonofont{SimSun}
\setCJKfamilyfont{song}[AutoFakeBold=2,BoldFont=SimHei]{SimSun}
\setCJKfamilyfont{hei}[AutoFakeBold=2,BoldFont=SimHei]{SimHei}
\setCJKfamilyfont{zhsong}[AutoFakeBold=2,BoldFont=SimHei]{SimSun}
\setCJKfamilyfont{zhhei}[AutoFakeBold=2,BoldFont=SimHei]{SimHei}
\setCJKfamilyfont{zhfs}[AutoFakeBold=2,BoldFont=SimHei]{FangSong}
\setCJKfamilyfont{callig}{STXingkai}
\newcommand{\calligfont}{\CJKfamily{callig}}
\XeTeXgenerateactualtext=1

% 全文保持黑色，不使用彩色链接或彩色表格
\hypersetup{colorlinks=false,hidelinks,linkcolor=black,citecolor=black,urlcolor=black,
  pdftitle={FinPlanRAG：面向金融混合查询的契约路由与锚点演化框架}}
% 预答辩/查重匿名版：不写入 pdfauthor，终稿恢复 pdfauthor={高深深}

% 表题在上，图题在下；图序/表序与题名之间空两个字符
\DeclareCaptionFont{sufeten}{\fontsize{10.5pt}{12.6pt}\selectfont}
\DeclareCaptionLabelSeparator{sufetwo}{\hspace{2em}}
\captionsetup[table]{position=top,font=sufeten,labelfont=bf,labelsep=sufetwo,skip=6pt,justification=centering}
\captionsetup[figure]{position=bottom,font=sufeten,labelfont=bf,labelsep=sufetwo,skip=6pt,justification=centering}
\renewcommand{\listfigurename}{图\quad 目\quad 录}
\renewcommand{\listtablename}{表\quad 目\quad 录}
\setlength{\heavyrulewidth}{1.0pt}
\setlength{\lightrulewidth}{0.5pt}
\renewcommand{\arraystretch}{1.15}

\ctexset{
  chapter/format={\centering\heiti\zihao{3}},
  chapter/beforeskip={24pt},
  chapter/afterskip={18pt},
  chapter/aftername={\qquad},
  chapter/afterindent=true,
  section/format={\centering\heiti\zihao{4}},
  section/name={第,节},
  section/number={\chinese{section}},
  section/afterindent=true,
  section/beforeskip={24pt},
  section/afterskip={6pt},
  section/aftername={\qquad},
  subsection/format={\heiti\fontsize{13pt}{15.6pt}\selectfont},
  subsection/name={,、},
  subsection/number={\chinese{subsection}},
  subsection/beforeskip={12pt},
  subsection/afterskip={6pt},
  subsection/aftername={},
  chapter/name={第,章},
  chapter/number={\chinese{chapter}},
  contentsname={目\quad 录},
  figurename={图},
  tablename={表}
}
\renewcommand{\thetable}{\thechapter.\arabic{table}}
\renewcommand{\thefigure}{\thechapter.\arabic{figure}}

\let\cleardoublepage\clearpage
\setlength{\parindent}{2em}
\setlength{\parskip}{0pt}
\makeatletter
\@afterindenttrue
\let\@afterindentfalse\@afterindenttrue
\makeatother
\linespread{1.37931}
\raggedbottom
\pagestyle{fancy}
\fancyhf{}
\fancyhead[C]{\zihao{5}\songti \nouppercase{\leftmark}}
\fancyfoot[C]{\zihao{5}\thepage}
\renewcommand{\headrulewidth}{0.4pt}
\renewcommand{\chaptermark}[1]{\markboth{#1}{#1}}
\renewcommand{\sectionmark}[1]{}

\newtheorem{proposition}{命题}[chapter]
\newtheorem{theorem}{定理}[chapter]
\newcommand{\E}{\mathbb{E}}
\newcommand{\indic}{\mathbf{1}}
\newcommand{\method}{FinPlanRAG}
\newcommand{\anchorevolve}{AnchorEvolve}
\newcommand{\sufehead}[1]{\chapter*{#1}\markboth{#1}{#1}\addcontentsline{toc}{chapter}{#1}}
\newcommand{\sufesec}[1]{\par\vspace{12pt}\noindent{\heiti\zihao{4}#1}\par\vspace{6pt}\phantomsection\addcontentsline{toc}{section}{#1}}

\title{\method：面向金融混合查询的契约路由与锚点演化框架}
% 预答辩/查重匿名版：不写入作者；终稿恢复 \author{高深深}
\author{}
\date{2026年11月}

\begin{document}
\setcounter{tocdepth}{3}
\setcounter{secnumdepth}{3}

% ==================== 封面（预答辩/查重匿名版） ====================
% 按学院通知：封面仅保留论文题目和专业，删除姓名、学号、学院、指导教师等个人信息；
% 原创性声明、致谢、个人简历在本阶段不排入，终稿阶段恢复。
\begin{titlepage}
\thispagestyle{empty}
\noindent
\begin{minipage}[t]{0.5\linewidth}
\vspace{0pt}
\includegraphics[width=2.6cm]{figures/sufe_logo.png}
\end{minipage}%
\begin{minipage}[t]{0.5\linewidth}
\raggedleft
{\zihao{5}\songti 学校代码：10272\par}
\end{minipage}

\vspace{1.1cm}
\begin{center}
{\calligfont\fontsize{40pt}{48pt}\selectfont 上海财经大学\par}
\vspace{0.45cm}
{\zihao{3}\bfseries SHANGHAI UNIVERSITY OF FINANCE AND\par}
{\zihao{3}\bfseries ECONOMICS\par}
\vspace{0.55cm}
{\zihao{2}\heiti 硕士学位论文\par}
\vspace{0.25cm}
{\zihao{3}\bfseries MASTER\quad DISSERTATION\par}
\end{center}

\vspace{2.6cm}
\begin{center}
{\zihao{4}\songti
\renewcommand{\arraystretch}{1.75}
\begin{tabular}{r@{\quad}l}
论文题目 & \underline{\makebox[9.6cm][c]{\textbf{FinPlanRAG：面向金融混合查询的}}} \tabularnewline
 & \underline{\makebox[9.6cm][c]{\textbf{契约路由与锚点演化框架}}} \tabularnewline
专\qquad 业 & \underline{\makebox[9.6cm][c]{金融工程与量化投资}} \tabularnewline
\end{tabular}\par}
\end{center}
\end{titlepage}

% ==================== 原创性声明与版权使用授权书（终稿恢复） ====================
% 按学院通知：预答辩与查重阶段不排入原创性声明和版权使用授权书，终稿阶段恢复。
\frontmatter
\pagestyle{fancy}
% \sufehead{学位论文原创性声明}
% {\zihao{-4}\songti
% \setlength{\parindent}{2em}
% \hspace*{2em}本人郑重声明：所呈交的学位论文，是本人在导师的指导下，独立进行研究工作所取得的成果。除文中已经注明引用的内容外，本论文不含任何其他个人或集体已经发表或撰写过的作品成果。对本文的研究做出重要贡献的个人和集体，均已在文中以明确方式标明。本人完全意识到本声明的法律结果由本人承担。\par
% \vspace{0.7cm}
% \noindent 论文作者签名：\includegraphics[height=1.0cm]{figures/signature.png} \hfill 日期：2026年11月10日\par
% }

% \vspace{0.8cm}
% \phantomsection
% \addcontentsline{toc}{chapter}{学位论文版权使用授权书}
% \begin{center}
% {\zihao{4}\heiti 学位论文版权使用授权书\par}
% \vspace{0.15cm}
% {\zihao{4}\heiti （硕士学位论文用）\par}
% \end{center}
% \vspace{0.15cm}
% {\zihao{-4}\songti
% \setlength{\parindent}{2em}
% \hspace*{2em}本人完全了解上海财经大学关于收集、保存、使用学位论文的规定，即：按照有关要求提交学位论文的印刷本和电子版本。上海财经大学有权保留并向国家有关部门或机构送交本论文的复印件和扫描件，允许论文被查阅和借阅。学校可以将学位论文的全部或部分内容编入有关数据库进行检索和传播，可以采用影印、缩印或扫描等复制手段保存、汇编学位论文。\par
% \vspace{0.7cm}
% \noindent 论文作者签名：\includegraphics[height=1.0cm]{figures/signature.png} \hfill 导师签名：\underline{\hspace{3cm}}\par
% \vspace{0.4cm}
% \noindent 日期：2026年11月10日 \hfill 日期：2026年11月10日\par
% }

% ==================== 中文摘要 ====================
\clearpage
\phantomsection
\addcontentsline{toc}{chapter}{摘\texorpdfstring{\quad}{ }要}
\begin{center}
{\zihao{-2}\heiti \method：面向金融混合查询的\\[0.2cm]
契约路由与锚点演化框架\par}
\vspace{0.2cm}
{\zihao{-2}\heiti 摘\quad 要\par}
\end{center}
\vspace{0.2cm}
\markboth{摘\quad 要}{摘\quad 要}
{\zihao{-4}\songti\sloppy
\setlength{\parindent}{0pt}
\setlength{\emergencystretch}{2em}
\hspace*{2em}金融文本分析同时包含确定信息问题与前瞻性断言。前者要求从截止时点以前的申报文件中读取或计算已陈述事实，后者要求判断尚未实现的未来结果。现有检索增强生成与断言验证通常预先固定输出接口，因而容易把前瞻性断言压缩为无保留的点答案，或把本可闭合的事实问题强制为分类裁决。本文提出FinPlanRAG，将输出契约表述为路由变量。LLM路由器依据查询与截止时点元数据选择义务检索或AnchorEvolve。义务检索以类型化义务组织证据获取，并由LLM阅读器返回带证据编号的闭合答案。AnchorEvolve抽取稳定关系、行业模式和实例三层锚点，依据可靠性、作用域和时效更新有限结果空间上的信念，保留可追溯的证据状态；校准读出再将该状态转换为点答案、可信集或弃权。\par
\hspace*{2em}本文在DocFinQA、227条XBRL核验收入指引断言和FinDVer上，使用DeepSeek Flash、GLM 5.3 Flash和Qwen3-8B，与六个具名基线进行统一评估。较大基座上的路由器恢复了接近oracle的混合效用，义务检索以较少证据取得更高准确率，AnchorEvolve的锚点后验在三套基座上均提供可测量的语义信号，包含LLM锚点信息的融合后验在GLM 5.3 Flash和Qwen3-8B上取得最高普通准确率；共享校准RF作为非LLM诊断，在决策损失上取得最低值。结果表明，契约路由负责选择推理接口，锚点演化负责组织前瞻性证据，校准读出负责控制最终风险。\par
\noindent{\zihao{-4}\heiti 关键词：}{\zihao{-4}\songti 金融问答\hspace{0.6em}检索增强生成\hspace{0.6em}前瞻性断言\hspace{0.6em}契约路由\hspace{0.6em}可信集与校准}
}

% ==================== 英文摘要 ====================
\clearpage
\phantomsection
\addcontentsline{toc}{chapter}{Abstract}
\begin{center}{\arial\bfseries\zihao{-2} Abstract\par}\end{center}
\vspace{0.6cm}
\markboth{Abstract}{Abstract}
{\arial\zihao{-4}\selectfont\sloppy
\setlength{\parindent}{0pt}
\setlength{\emergencystretch}{3em}
\hspace*{2em}Financial text analysis combines determinate questions, whose answers are stated or computable at a cutoff, with forward-looking claims, whose outcomes have not yet been realized. Existing retrieval-augmented generation and claim-verification systems usually fix the output interface in advance, which can reduce a forward-looking claim to an unqualified point answer or force a closed fact question into a classification verdict. We introduce FinPlanRAG, which treats the output contract as a routing variable. An LLM router selects obligation retrieval or AnchorEvolve from the query and cutoff metadata. Obligation retrieval organizes evidence acquisition around typed obligations and returns a closed answer with evidence identifiers. AnchorEvolve extracts stable-relation, industry-schema, and instance anchors, qualifies their reliability, scope, and timeliness, and updates a belief over a finite outcome space while preserving an auditable evidence state. A calibrated readout converts this state into a point answer, a credible set, or an abstention.\par
\hspace*{2em}We evaluate FinPlanRAG on DocFinQA, 227 XBRL-verified revenue-guidance claims, and FinDVer with DeepSeek Flash, GLM 5.3 Flash, and Qwen3-8B against six named baselines. The routers built on the two larger backbones recover near-oracle mixed utility, obligation retrieval improves accuracy with fewer evidence passages, and the anchor posterior provides measurable semantic signal across all three backbones. The LLM-dependent fused posterior attains the highest standard accuracy on GLM 5.3 Flash and Qwen3-8B. A shared calibrated RF is reported separately as a non-LLM diagnostic and yields the lowest decision loss. These results establish a division of labor in which contract routing selects the inference interface, anchor evolution organizes forward-looking evidence, and calibrated readout controls output risk. FinPlanRAG therefore provides an auditable and calibratable design for financial text workloads that combine factual retrieval with future judgment.\par
}
\vspace{0.4cm}
{\sloppy\noindent{\arial\bfseries\zihao{-4} Key Words:} {\arial\zihao{-4}\selectfont financial QA\hspace{0.8em}RAG\hspace{0.8em}forward-looking claims\hspace{0.8em}contract routing\hspace{0.8em}calibrated credible sets}\par}

% ==================== 目录 ====================
\clearpage
\tableofcontents
\markboth{目\quad 录}{目\quad 录}
\clearpage
\phantomsection
\addcontentsline{toc}{chapter}{图\texorpdfstring{\quad}{ }目\texorpdfstring{\quad}{ }录}
\listoffigures
\markboth{图\quad 目\quad 录}{图\quad 目\quad 录}
\clearpage
\phantomsection
\addcontentsline{toc}{chapter}{表\texorpdfstring{\quad}{ }目\texorpdfstring{\quad}{ }录}
\listoftables
\markboth{表\quad 目\quad 录}{表\quad 目\quad 录}
\clearpage

\mainmatter
\zihao{-4}
\setlength{\parindent}{2em}
\zihao{-4}

% ==================== 第1章 ====================
\chapter{绪论}
\label{ch:intro}

\section{研究背景与问题提出}

金融文本分析连接财务报告理解、事实检索与业绩判断，但真实查询并不服从单一的问答契约。确定信息问题要求系统从截止时点以前的10-K、10-Q、8-K及其表格、脚注和XBRL事实中读取或计算已经陈述的结果，评价重点是答案准确性与证据可定位性；前瞻性断言则涉及尚未实现的营收指引、产能投放、监管审批或合同履约，其真值在截止时点之后才会揭示，合理输出可能是点判断、可信集或弃权。管理层预测研究也表明，指引同时包含区间、时点和可验证性，不能等同于已实现业绩\citep{hirst2008management}。围绕长文档、表格推理和证据获取，金融问答与检索增强生成已经形成系统方法\citep{lewis2020rag,chen2021finqa,zhu2021tatqa,islam2023financebench,reddy2024docfinqa}，主动检索、自适应检索和智能体框架进一步优化了检索时机与证据使用\citep{jiang2023flare,jeong2024adaptiverag,yao2023react,shinn2023reflexion}；金融断言验证研究将陈述映射为支持、反驳或证据不足\citep{zhao2024findver,rangapur2025finfact}，选择性预测则通过拒识控制低置信度输出\citep{chow1970optimum,geifman2017selective}。这些研究分别改进了证据获取、断言核验和风险控制，却通常预先固定回答接口。固定契约用于混合负载时，确定问题可能被压缩为分类裁决，前瞻性断言也可能被阅读器输出为无保留的点答案。尚未被充分处理的关键问题因此不是单一检索器或分类器的性能，而是系统能否先判断一条查询应遵循事实检索契约还是前瞻性判断契约，并在相应契约下维护证据状态。本文据此提出FinPlanRAG：给定查询$q$与截止时点$c$，LLM路由器选择义务检索或AnchorEvolve；前者以类型化义务驱动时点取证并返回带证据编号的闭合答案，后者以稳定关系、行业模式和实例三层锚点组织前瞻性证据，依据可靠性、作用域和时效更新结果信念。点答案、可信集或弃权由校准读出生成，使契约选择、证据演化与风险控制在同一时点约束下形成可审计链条。

\section{研究内容与方法}

本文按照查询契约识别、分支证据建模和统一实证检验展开。首先，将金融混合查询形式化为确定信息问题与前瞻性断言问题的并集，并以截止时点界定可用信息，所有检索、特征构造、校准和读出选择均遵守时点约束，以控制前视偏差\citep{lopezdeprado2018advances}。其次，构建由LLM路由器、义务检索和AnchorEvolve组成的总体框架。路由器仅依据查询与截止时点元数据返回契约标签、置信度和理由；确定信息问题被解析为类型化义务，系统围绕未覆盖义务迭代检索，由LLM阅读器判断证据闭包并生成带编号答案；前瞻性断言被映射到有限结果空间，AnchorEvolve抽取L0稳定关系、L1行业模式和L2实例锚点，通过五类资格判定更新信念，校准读出据此生成点答案、可信集或弃权。最后，本文在DocFinQA、XBRL核验收入指引和FinDVer上，使用DeepSeek Flash、GLM 5.3 Flash和Qwen3-8B进行同基座比较，并分别检验路由效用、答案质量、证据成本、锚点增量信息、概率校准和非对称决策损失。这样的研究设计使契约识别、证据组织与风险控制能够被分层验证，同时保留端到端系统在混合负载上的整体评价。

\section{主要贡献}

本文的主要贡献形成于查询层、证据层和评价层的连续设计。查询层将回答形式选择表述为混合负载上的契约路由，把“事实在截止时点是否已经存在”置于检索与生成之前，使闭合问答与前瞻性裁决可以共享时点约束而保持不同输出接口。证据层提出AnchorEvolve的金融断言实例化，以L0稳定关系、L1行业模式和L2实例锚点记录证据的取值、作用域、来源、时间与可靠性，通过五类资格判定形成可更新、可追溯的信念状态，并让接受的证据按类型和公开揭示日更新Beta可靠性；类型库扩展作为验证集门控的训练步骤，冻结后才用于测试。义务检索以类型化义务和阅读器充分性控制确定问题的证据获取。决策层用全动作损失感知读出把点答案、非单元素裁决集和弃权放在同一个期望损失下比较。评价层在DocFinQA、227条收入指引断言、470条A股年报断言和150条FinDVer断言上，以三套基座和同基座基线检验路由、义务检索、锚点推断、后验融合、类型库扩展、数值桥和外部迁移。由此，FinPlanRAG不仅给出端到端结果，也保留证据编号、锚点来源、资格判定和信念变化，为金融分析中的复核、问责与事后审计提供可复用接口。

% ==================== 第2章 ====================
\chapter{文献综述}
\label{ch:related}

本章围绕金融问答与断言验证、检索增强生成与自适应检索、选择性预测与可信集、锚点约束与时点证据四条线索梳理相关研究，并在此基础上归纳研究缺口与关键挑战。

\noindent\textbf{一、金融问答与断言验证}

金融问答基准测试模型在申报文件、表格和报告上的数值推理能力。FinQA研究基于财务报表的数值推理\citep{chen2021finqa}，TAT-QA研究表格与文本的混合推理\citep{zhu2021tatqa}，FinanceBench研究真实金融文档上的开放问答\citep{islam2023financebench}，FinBen研究金融领域的综合评测\citep{xie2024finben}。DocFinQA在FinQA问题的基础上加入完整10-K上下文，把证据选择从段落级扩展到长文档级\citep{reddy2024docfinqa}。这些基准的共同点是答案在文档中已经存在，评价指标主要是精确匹配或词级F1。它们衡量的是阅读器能否在给定文档中找到并计算答案，而不是系统能否判断查询是否应当以闭合答案回应。

金融申报文件的结构与通用问答语料不同。10-K和10-Q通常包含多级标题、跨页表格、脚注和附注，同一指标可能在正文、表格和管理层讨论中以不同口径出现；8-K则按事项编号披露业绩、指引、并购和监管事件。XBRL数据提供结构化财务事实，但季度口径、单位、重述和可用日期仍需与文本对齐。通用阅读器擅长语义匹配，却不一定能处理期间、单位和证据闭合，因此金融问答需要把文档结构、时点约束和数值计算纳入同一证据策略。

断言验证基准测试模型判断一条陈述是否被证据支持、反驳或无法判断。FinDVer面向金融领域的断言验证，要求模型在长文档和混合内容中给出裁决和解释\citep{zhao2024findver}；Fin-Fact面向证据驱动的金融事实核查，强调多模态证据和解释生成\citep{rangapur2025finfact}。这类任务的输出是离散裁决，评价指标是裁决准确率和解释质量。验证器通常假设陈述已经被给定，且证据已经被检索，因而没有回答“这条陈述是否应当进入验证流程”的问题。

近期金融RAG系统进一步加入类型化证据、期间约束、来源追踪和语料对齐检索\citep{mansouri2026finrank,jiang2026finrate,chen2026hcrag,chi2026finsagent,zhou2026fincards}。这些系统在证据表示和检索配方上明显优于通用RAG，但输出接口仍通常是固定答案或固定裁决。它们改善了到达阅读器的证据，却没有改变阅读器与验证器之间的契约边界。对于同时包含事实问题和前瞻性断言的金融分析场景，这种固定接口仍然是主要限制。

\noindent\textbf{二、检索增强生成与自适应检索}

RAG将检索器与生成器组合，假设答案可以由检索到的上下文拼装\citep{lewis2020rag}。主动检索和自适应检索进一步决定何时检索或检索多少上下文\citep{jiang2023flare,jeong2024adaptiverag}，智能体框架则交错进行推理、检索和自我修正\citep{yao2023react,shinn2023reflexion}。Self-RAG通过自我反思决定是否检索和如何评价证据\citep{asai2024selfrag}，CRAG在检索质量不足时触发纠正性检索\citep{yan2024crag}，FLARE在生成过程中按置信度触发检索\citep{jiang2023flare}，Adaptive-RAG按问题复杂度选择检索策略\citep{jeong2024adaptiverag}。

这些方法改善了检索过程和证据使用，但输出通常仍是生成的答案。它们可以决定“是否检索”和“检索多少”，却不能决定“返回答案、可信集还是弃权”。在确定信息问题上，这种输出是合适的；在前瞻性断言上，单一生成答案无法表达结果分布、证据作用域和决策风险。检索策略的自适应并不等于输出契约的自适应，二者是两个不同层次的决策。

本文的义务检索分支借用了自适应检索的迭代思想，但把停止条件从“检索轮数”改为“义务覆盖与阅读器充分性”。这一改动使证据预算服务于问题的语义结构，而不是服务于固定的$k$。AnchorEvolve则把输出契约从生成答案扩展为点答案、可信集和弃权的组合，使系统可以在证据不足时保留不确定性。

\noindent\textbf{三、选择性预测、校准与可信集}

拒识选项在识别错误和拒识率之间进行权衡\citep{chow1970optimum}，选择性分类通过置信度阈值在低置信度输入上弃权\citep{geifman2017selective}。现代分类器通常存在校准误差，温度缩放是常用基线\citep{guo2017calibration}；代价敏感学习为不同错误赋予不同价格\citep{elkan2001cost}；分裂共形预测则把得分转换为在可交换性下具有分布无关覆盖率的集合\citep{vovk2005algorithmic,angelopoulos2021gentle}。分类共形方法进一步构造自适应集合\citep{romano2020classification}，分布漂移下的自适应共形方法则允许覆盖率随时间调整\citep{gibbs2021adaptive}。

FinAbstain进一步把时点检索、多模态证据、证据核验、不确定性校准和显式弃权组合成选择性金融预测流程\citep{torres2026finabstain}。它是与本文模糊分支最接近的工作，但仍然把查询预先视为预测任务，没有在闭合问答与前瞻性裁决之间做契约路由；其证据核验也是不确定性聚合之前的过滤，而不是独立的多类别资格判定阶段。

这些方法提供了损失、置信度和集合构造的工具，但通常假设任务和回答接口已经固定。选择性分类器在“答案”和“弃权”之间选择，验证器在“支持”“反驳”和“证据不足”之间选择，共形方法在给定标签空间上构造集合。它们没有解决“标签空间本身是否适用于当前查询”的问题。对于前瞻性断言，标签空间是结果桶，证据是锚点，输出还需要说明证据的作用域和时效；这些结构不在标准选择性预测的范围内。

本文借用选择性预测的损失框架和共形集合的覆盖思想，但把标签空间、证据资格和输出动作放在同一个锚点状态上。校准随机森林和温度缩放用于控制后验概率，损失感知门用于在答案、可信集和弃权之间选择。与标准共形方法相比，本文的可信集由后验质量构造，覆盖率在实验中被报告，但并未假设交换性；这是金融时点数据与标准共形设定之间的差异。

\noindent\textbf{四、锚点、部分识别与时点证据}

部分识别表示与现有证据一致的结果集合\citep{manski2003partial}。在金融机器学习中，模型决策只能使用当时可得的信息，否则会产生前视偏差\citep{lopezdeprado2018advances}。本文借用约束和可行集的思想，将锚点定义为来源与假设之间的关系：硬锚点可以排除某个结果桶，软锚点只重新加权结果信念；证据过滤和模式过滤均按截止时点执行。

与部分识别相比，本文的锚点不是完整的结构模型，而是由LLM从文本中抽取的语义约束。锚点带有类型、取值、作用域、来源、时间、可靠性和有效期，因而可以被逐条检查。部分识别给出的是在给定模型下与证据一致的参数集合，锚点给出的是在给定断言下与证据一致的结果桶约束。两者都强调“证据不足时不要强行点识别”，但本文更关注证据的可审计性和可更新性。

时点约束贯穿本文的两个分支。义务检索只接受可用日期不晚于截止时点的文档，AnchorEvolve只接受截止时点之前的证据，随机森林只使用截止时点之前可得的特征，校准和权重选择只在开发集与验证集上进行。这一约束使“当时可得信息”成为系统设计的一部分，而不是事后评价的附加条件。

\noindent\textbf{五、研究缺口与关键挑战}

现有研究分别推进了长文档检索、金融问答、断言验证和低置信度拒识，但尚未把查询契约选择与前瞻性证据演化纳入同一框架。本文认为，这一缺口可以分解为三个相互衔接的挑战。三个挑战分别对应查询层、证据层和决策层，也构成本文框架的设计依据。

第一个挑战是查询契约的识别。金融查询的表面形式与真实契约并不一一对应。“预期股利”“未来支付金额”等表述包含未来时态，但答案可能在文件中已经给出；管理层指引包含精确数字，却是关于尚未实现结果的前瞻性断言。系统如果按关键词或时态路由，会把确定问题送入裁决分支，或把前瞻性断言送入阅读器。真正的判断依据是查询所要求的事实是否在截止时点已经存在，以及输出应当采用闭合答案还是可校准判断。这一挑战的难点不在分类器本身，而在于契约是查询、证据与时点之间的关系，而不是查询文本的主题标签。

第二个挑战是前瞻性证据的状态化。现有验证器通常把断言直接映射为支持、反驳或证据不足，阅读器则直接生成答案。两者都没有保留证据的作用域、可靠性、时效和来源，也没有在证据变化时更新判断。对于前瞻性断言，证据往往只支持某个结果桶、只在某个时间范围内有效、只对某个公司或行业成立；把这些信息压缩为一个标签，系统就无法解释判断为何改变，也无法在结果揭示后更新可靠性。本文认为，前瞻性断言需要一个可审计的中间状态，使证据到信念、信念到决策的每一步都可以追溯。

第三个挑战是语义判断与风险校准的分工。LLM能够从文本中抽取锚点、解释证据和形成语义判断，但其概率输出往往存在类别偏置和校准误差；统计模型可以使用结构化特征进行概率校准，却缺少对文本证据的理解。若让LLM直接输出最终裁决，系统可能过度自信；若只用统计模型，系统又无法利用文本中的语义约束。本文认为，合理的分工是让LLM构造语义状态，让校准程序控制决策风险，并在验证集上选择读出方式，而不是在测试集上比较后择优。

上述三个挑战对应本文的三项贡献。契约识别由查询契约路由回应；证据状态化由AnchorEvolve的三层锚点、五类资格判定和信念更新，以及义务检索回应；语义与校准的分工由校准后端、损失感知门和可信集读出回应。三项贡献分别对应查询层、证据层和决策层，构成逐层回应。

% ==================== 第3章 ====================
\chapter{研究设计：研究问题、假设与总体框架}
\label{ch:framework}

\section{研究问题}

金融文本分析中的查询可以分为两类。确定信息问题要求返回申报文件在截止时点已经陈述或可以计算的事实，答案通常是数值、段落或明确的预测。前瞻性断言问题要求判断尚未实现的未来结果，其真值取决于截止时点之后披露的信息，正确输出可能是一个点答案、一个可信集，或者在证据不足时弃权。设查询为$q$，截止时点为$c$，工作负载是确定问题集合$\mathcal Q_D$与前瞻性断言集合$\mathcal Q_V$的混合。系统调用LLM路由器
\begin{equation}
r(q,c)\in\{D,V\}
\end{equation}
判断查询属于哪一类，再将确定查询送入义务检索分支，将前瞻性断言送入锚点演化分支。

两个分支返回不同的输出形式。义务检索分支返回闭合答案$a$、置信度与证据编号，证据不足时返回\texttt{insufficient}。AnchorEvolve返回点答案$h^*$、可信集$R_\alpha$、动作$\pi$与信念$b$，动作可以是回答、可信集裁决或弃权。路由器只读取原始查询和截止时点元数据，不读取“这是断言”的任务标签，也不读取测试结果。

上述问题可以展开为三个研究问题：
\begin{enumerate}[label=RQ\arabic*.,leftmargin=2em,itemsep=2pt,topsep=2pt]
\item 契约识别：仅依据查询与截止时点元数据，系统能否在混合负载上区分确定信息契约与前瞻性断言契约？
\item 确定信息取证：类型化义务与覆盖驱动的证据获取，能否在更低证据成本下提高确定信息问答的准确率与覆盖率？
\item 前瞻性裁决：三层锚点状态与校准读出，能否在保持分类准确率的同时降低决策损失，并提供可信集或弃权？
\end{enumerate}
三个问题分别对应查询层、证据层和决策层。RQ1检验路由器能否识别契约；RQ2检验义务检索能否以更低成本提高确定问答的准确率和覆盖率；RQ3检验锚点状态与校准读出能否把语义信息转换为低风险决策。

\section{研究假设}

三个研究问题对应三个可检验假设。假设的表述同时给出检验口径和否证条件，使后文实验能够直接检验框架是否有效，而不是只描述框架结构。

\textbf{H1（契约可识别性）}：仅依据查询与截止时点元数据，LLM路由器能够在混合负载上区分确定信息契约与前瞻性断言契约；在较强基座上，路由混合效用接近oracle，并高于Always-RAG和Always-Anchor。检验口径为573条混合负载上的路由准确率、混淆矩阵和混合效用。若路由器在多数基座上低于较好的固定分支，则H1不成立。

\textbf{H2（义务取证的效率）}：类型化义务与覆盖驱动的证据获取，能够在低于固定16段的平均证据成本下，提高DocFinQA的精确匹配和答案覆盖率，并迁移到FinDVer。检验口径为DocFinQA的300题固定样本和FinDVer testmini上的准确率、宏F1、平均段落数和覆盖率。若分支在多数基座上不高于固定16段阅读器，或证据成本没有下降，则H2不成立。

\textbf{H3（语义状态与校准读出的分工）}：三层锚点状态能够提供结构化语义信号，校准后端与损失感知门能够把该状态转换为低风险决策；最终系统在保持普通准确率的同时降低非对称决策损失，并在锚点消融和输出形式消融中表现出可解释的差异。检验口径为227条XBRL核验断言上的准确率、平衡准确率、宏F1、ECE、Brier、决策损失、可信集覆盖率和平均集合大小。若最终系统在所有基座上的决策损失不低于最强具名基线，或去除锚点后宏F1没有下降，则H3不成立。

三个假设共同构成对研究问题的可否证回答。H1对应查询层，H2对应确定信息的证据层，H3对应前瞻性断言的证据层与决策层。第\ref{ch:empirical}章按这三个假设组织实验，而不按数据集顺序堆叠结果。

\section{总体框架}

FinPlanRAG由三个模块组成。路由器是一个独立的LLM调用，负责判断查询契约；义务检索分支负责确定信息问题；AnchorEvolve分支负责前瞻性断言。路由器与两个分支之间通过结构化输出连接，路由器不读取分支输出，分支不改变路由结果。这样的结构使路由错误和分支错误可以被分开诊断：路由错误表现为查询被送入不匹配的契约，分支错误表现为在正确契约下的证据或判断错误。

路由器读取查询文本与截止时点元数据，返回路由标签、置信度和简短理由。提示词定义两个契约：确定信息请求要求返回申报文件在截止时点陈述或可计算的事实、数值、段落、估计或预测；前瞻性断言要求判断尚未实现的未来结果，例如实际营收是否落在指引区间内、产品是否按时发布、目标是否达成。提示词明确说明管理层指引是前瞻性断言，不是已实现事实，即使其中包含数字。

路由器不读取已实现结果，也不读取两个分支的输出。它只看到查询和截止时点元数据，因此在部署时可以按查询缓存，成本是一次短LLM调用。路由器的错误类型也具有语义含义：确定问题被送到AnchorEvolve，通常是因为查询文本包含未来时态或预期措辞；前瞻性断言被送到义务检索，通常是因为断言中包含可检索的数字或期间。这些错误为改进路由器提供了可解释的信号。

两个分支共享路由器，但不共享输出契约。义务检索分支返回闭合答案，AnchorEvolve返回点答案、可信集或弃权。共享路由器使系统只需一次契约判断；不共享输出契约使每个分支可以针对自己的证据结构和风险结构设计。图\ref{fig:framework}给出总体结构，第\ref{ch:determinate}章和第\ref{ch:anchor}章分别给出两个分支的完整设计。

\begin{figure}[H]
\centering
\resizebox{0.98\linewidth}{!}{%
\begin{tikzpicture}[
  font=\normalsize,
  >=Latex,
  box/.style={draw, rounded corners=3pt, align=center, text width=2.6cm, minimum height=0.9cm, inner sep=3pt, fill=white},
  panel/.style={draw, rounded corners=6pt, inner sep=7pt, fill=gray!5},
  arrow/.style={->, thick, draw=black!65},
  lbl/.style={font=\footnotesize, text=black!70}
]
\node[box, fill=gray!10] (q) at (0,0) {查询 $q$\\截止时点 $c$};
\node[box, fill=gray!18, text width=2.6cm] (r) at (2.7,0) {LLM 路由器\\契约判断};

\node[box] (o1) at (5.6,2.0) {类型化义务};
\node[box] (o2) at (8.6,2.0) {时点检索};
\node[box] (o3) at (11.6,2.0) {LLM 阅读器};

\node[box] (a1) at (5.6,-2.0) {三层锚点};
\node[box] (a2) at (8.6,-2.0) {资格与信念};
\node[box] (a3) at (11.6,-2.0) {校准与损失门};

\begin{scope}[on background layer]
\node[panel, fit=(o1)(o3), label={[font=\bfseries]above:义务检索分支}] {};
\node[panel, fit=(a1)(a3), label={[font=\bfseries]below:AnchorEvolve 分支}] {};
\end{scope}

\node[box, text width=2.1cm] (ans) at (14.6,2.0) {闭合答案};
\node[box, text width=2.1cm] (pt) at (14.6,0.6) {点答案};
\node[box, text width=2.1cm] (cs) at (14.6,-0.6) {可信集};
\node[box, text width=2.1cm] (ab) at (14.6,-2.0) {弃权};
\node[lbl, above] at (14.6,2.65) {输出};

\draw[arrow] (q) -- (r);
\draw[arrow] (r.east) -- ++(0.7,0) |- (o1.west);
\draw[arrow] (r.east) -- ++(0.7,0) |- (a1.west);
\draw[arrow] (o1) -- (o2);
\draw[arrow] (o2) -- (o3);
\draw[arrow] (a1) -- (a2);
\draw[arrow] (a2) -- (a3);
\draw[arrow] (o3) -- (ans);
\draw[arrow] (a3) -- (pt);
\draw[arrow] (a3) -- (cs);
\draw[arrow] (a3) -- (ab);

\draw[dashed, draw=black!45] (0.9,3.0) -- (0.9,-3.0);
\node[lbl, above] at (0.9,3.05) {时点约束 $V_c$};
\end{tikzpicture}
}
\caption{FinPlanRAG总体框架。路由器把查询送入义务检索分支或AnchorEvolve；前者返回闭合答案，后者返回点答案、可信集或弃权。左侧虚线表示时点约束。}
\label{fig:framework}
\end{figure}


\section{时点约束与评价协议}

系统在截止时点$c$形成判断，可用信息由可采纳证据集合$V_c$给出。设证据$e$的可用时间为$a(e)$，则
\begin{equation}
V_c=\{e:a(e)\le c,\ e\text{在}c\text{时未被撤回}\}.
\end{equation}
策略$\pi$在$c$时有效，当且仅当它是$\sigma(V_c)$可测的。这一条件排除了使用未来结果、未来披露或未来修订的证据。路由器不读取结果，锚点抽取器不接收结果，随机森林只使用截止时点之前可得的特征，校准和权重选择只在开发集与验证集上进行。测试结果不参与任何拟合或选择。

评价口径与输出契约对应。确定信息问题使用精确匹配、词级F1、答案覆盖率和证据成本；前瞻性断言使用准确率、平衡准确率、宏F1、MCC、AUROC、AUPRC、ECE、Brier、决策损失、可信集覆盖率和平均集合大小；混合负载使用路由准确率、混淆矩阵和混合效用。所有任务共享同一时点约束和同一套测试集不参与选择的原则，因此不同分支的结果可以在同一框架下比较。

评价协议由三部分组成。第一，主结果使用配对自助法给出置信区间，准确率差异使用McNemar检验；所有区间在断言或问题级别重采样。第二，随机森林使用五个折种子，测试年份的训练集来自其他公司折，且结果日期早于测试年份1月1日；测试公司不出现在训练集中。第三，所有LLM响应按请求负载的哈希缓存，提示词、运行脚本、逐条预测和表格生成脚本随论文发布，API密钥只从环境变量读取。

这一协议的目的是让每个假设都有对应的证据边界。H1对应573条混合负载；H2对应DocFinQA的300题固定样本和FinDVer testmini；H3对应227条XBRL核验断言、锚点消融和输出形式消融。

% ==================== 第4章 ====================
\chapter{义务检索分支：Obligation Retrieval}
\label{ch:determinate}

义务检索分支可以概括为四步：把问题表示为类型化义务，构建时点证据空间，按未覆盖义务迭代取证，由阅读器生成带证据编号的闭合答案。四步依次对应本章四节。

\section{状态表示：类型化义务}

义务检索分支把自然语言问题绑定为可检索、可核对的类型化义务。义务元组定义为
\begin{equation}
o=(e,f,y,p,m,u),
\end{equation}
其中$e$为实体，$f$为表单或文档类型，$y$为财年，$p$为财期，$m$为指标，$u$为单位。以“某公司2023财年第四季度营业收入是多少”为例，实体是公司，表单是10-K或10-Q，财年是2023，财期是第四季度，指标是营业收入，单位是美元。义务项从问题和查询元数据中抽取；抽取器只负责识别约束，不负责回答。

文档只有在可用日期不晚于截止时点、且元数据与义务一致时才被接受。元数据一致性检查覆盖实体、表单、财年、财期、指标和单位。与只按问题文本检索相比，义务元组把实体、期间和单位变成可检查的约束。例如，问题文本中的“2023年”可能指财年、自然年或季度，BM25无法区分；义务元组把“2023财年第四季度”固定为可核对的期间，从而减少跨期文档和单位不一致带来的错误。

类型化义务还提供了证据覆盖的度量。设问题$q$的义务项集合为$T(q)$，已获取证据集合为$A_t$，则未覆盖义务项为
\begin{equation}
U_t=T(q)\setminus\{t\in T(q):\exists a\in A_t,\ t\in a\}.
\end{equation}
未覆盖项既用于构造下一轮检索查询，也用于判断何时停止。义务覆盖不是答案正确性的充分条件，但它把“证据是否覆盖了问题的关键约束”变成可观察的状态，使系统可以在证据不足时继续检索或返回证据不足，而不是在固定轮数后强行生成答案。

\section{证据空间：时点语料与块索引}

分支把文档切分为段落块和表格块。段落块按行合并，单块不超过1200字符；表格块按连续表格行合并，单块不超过4000字符。每个块记录文档编号、可用时间、文本、词项和元数据。切块的目标是保留表格行与表头的对应关系，同时避免单块过长导致阅读器输入超出预算。对于跨页表格，分支保留表头并在块内重复列名，以减少阅读器对列含义的误判。

分支使用BM25建立索引，按问题与未覆盖义务项检索候选块。BM25的优势是无需训练、可复现、对长文档稳定；缺点是不理解语义，因此需要义务覆盖和阅读器反馈补足。检索查询由原问题与$U_t$拼接而成，候选块按得分排序，去重后最多新增两个块。这样做的原因是：一次加入过多块会稀释阅读器的注意力，也会增加token成本；一次加入过少块则可能需要更多轮检索。两块的设置是在证据覆盖和调用成本之间的折中。

时点过滤在检索前执行。只有可用日期不晚于截止时点的文档进入索引；被后续版本撤回或替代的文档被排除。对于同一文档的多个版本，分支保留截止时点之前的最新版本。时点过滤使系统不会使用未来披露的信息，也避免了“文档存在但当时尚未发布”的前视偏差。

\section{取证策略：覆盖驱动的迭代检索}

第$t$轮检索后，分支把新增块加入证据集$A_t$，调用LLM阅读器判断义务是否闭合。停止条件有三类：阅读器报告证据充分且$U_t=\varnothing$；阅读器给出非空答案、置信度不低于0.8且$U_t=\varnothing$；预算$B$耗尽。默认$B=8$，并报告$B=12$和$B=16$。前两类条件把停止权交给证据状态和阅读器判断，第三类条件提供成本上限。

覆盖驱动的取证把停止变量从检索前的固定预算移到义务覆盖和阅读器充分性上。对于简单问题，分支可以在两三轮内闭合；对于复杂问题，分支可以继续检索直到预算耗尽。这种自适应性使平均证据成本低于固定预算，同时在难题上保留足够的证据预算。

置信度门控是覆盖驱动取证的补充。当阅读器置信度低于0.8或报告证据不足时，分支回退到固定16段证据集重新阅读。回退不是对每个问题支付16段成本，而是对规划器无法闭合的问题支付更高成本。实验表明，回退只在规划器无法闭合的问题上触发，DeepSeek Flash上约为24\%，GLM 5.3 Flash上比例更高；它弥补了规划器单独的准确率缺口，同时保持平均段落数低于固定16段。完整流程见算法\ref{alg:obligation}。



\section{答案生成：阅读器与闭包}

阅读器接收问题与已获取块，返回结构化答案
\begin{equation}
(\text{answer},\text{confidence},\text{sufficient},\text{missing},\text{evidence\_ids},\text{rationale}).
\end{equation}
答案要求为短字符串；若证据不足，返回\texttt{insufficient}并将\texttt{sufficient}置为\texttt{false}。分支把阅读器的\texttt{missing}字段拼入下一轮查询，形成阅读器驱动的取证。例如，阅读器报告缺少“第四季度”或“单位”，下一轮查询会加入这些词项，从而把检索引向缺失的证据。

闭包规则要求答案附带证据编号。每个答案都可以追溯到具体段落或表格块，证据不足时有明确的缺失信息。这一规则使义务检索分支的输出可以被下游审计：读者可以检查答案是否由证据支持，也可以检查系统在证据不足时是否选择弃权而不是猜测。

解析失败在评价中计为错误。阅读器被要求返回JSON，但基座可能返回解释性文字或格式错误的JSON。分支对解析失败不进行静默修复，而是记录为错误，以避免评价口径被解析器放宽。Qwen3-8B在DocFinQA上返回较多解释性文字，导致精确匹配接近地板；这一结果作为基座能力边界报告，不用于义务检索分支的主比较。

\begin{algorithm}[H]
\caption{义务检索的单轮取证}
\label{alg:obligation}
\begin{algorithmic}[1]
\Require 问题$q$，截止时点$c$，义务项$T(q)$，预算$B$
\Ensure 答案$a$或\texttt{insufficient}
\State $A\gets\varnothing$; $U\gets T(q)$
\For{$t=1,\dots,B$}
  \State $C_t\gets\textsc{BM25}(q\cup U,\,V_c)$
  \State $A\gets A\cup\textsc{Top2}(C_t\setminus A)$
  \State $(a,\text{conf},\text{suff},\text{miss},\text{ids})\gets\textsc{Reader}(q,A)$
  \State $U\gets T(q)\setminus\textsc{Covered}(A)$
  \If{$\text{suff}=\text{true}$ and $U=\varnothing$}
    \State \Return $a$
  \EndIf
  \If{$\text{conf}\ge0.8$ and $a\ne\varnothing$ and $U=\varnothing$}
    \State \Return $a$
  \EndIf
  \State $q\gets q\cup\text{miss}$
\EndFor
\State $A\gets\textsc{Fixed16}(q,V_c)$
\State $(a,\text{conf},\text{suff},\dots)\gets\textsc{Reader}(q,A)$
\If{$\text{suff}=\text{true}$}
  \State \Return $a$
\Else
  \State \Return \texttt{insufficient}
\EndIf
\end{algorithmic}
\end{algorithm}

义务检索的时间成本主要由BM25检索和阅读器调用组成。BM25在本地索引上完成，不产生API成本；阅读器调用次数由停止条件决定，最坏情况下为预算$B$加上一次固定16段回退。设文档块数为$N$，单次BM25检索的复杂度为$O(N\log N)$，相对于LLM调用可以忽略。

可审计性来自证据编号和缺失信息。每个答案都附带实际使用的段落或表格块编号，证据不足时返回具体的缺失项。审计者可以检查答案是否由证据支持，也可以检查系统在证据不足时是否选择弃权而不是猜测。

% ==================== 第5章 ====================
\chapter{AnchorEvolve：前瞻性断言分支}
\label{ch:anchor}

AnchorEvolve是FinPlanRAG的前瞻性断言分支。本章先给出状态表示、证据来源和总体结构，再给出完整算法，包括证据构造、锚点抽取、五类资格判定、时间演化、类型库扩展、信念更新、全动作读出、数值桥和披露日期断言。

% ==================== 第4章 ====================

\section{设计原则}

FinPlanRAG的设计遵循六条原则，其核心模块AnchorEvolve承担证据状态化、资格判定、时间演化和信念更新。这些原则不是附加的功能列表，而是对第\ref{ch:framework}章所述研究缺口的直接回应。

第一，契约先行。系统在检索证据之前先确定公司、目标期间、指标、结果空间和截止时点。契约把“回答什么”写成显式对象，使不同方法可以在同一口径下比较，也使用户可以在运行前检查阈值是否合理。契约不是系统的内部细节，而是审计报告的第一部分。

第二，来源落地。每条锚点必须引用一个具体的证据来源，并给出逐字引文。系统不接受没有来源的判断，也不接受只在理由中描述证据而不给出引文的锚点。来源落地使锚点可以被逐条复核，也使错误的锚点可以被定位到具体的网页或段落。

第三，资格先于权重。锚点在进入信念更新之前先经过资格判定。空泛、过期、对抗三类锚点不进入更新；硬锚点和软锚点分别处理。这个顺序很重要：如果先把所有锚点加权再过滤，被拒绝的锚点仍然会影响后验；如果先过滤再加权，每个被采纳的锚点都有明确的资格理由。

第四，揭示后演化。锚点类型的可靠性只在结果公开揭示之后更新，且只使用揭示日早于当前截止时点的结果。这个约束使可靠性演化在回测和部署中具有相同的含义，避免把未来结果泄漏到时点判断。

第五，损失感知。输出不是固定的点答案，而是在点答案、非单元素裁决集和弃权三个动作上最小化期望损失。损失函数把过度乐观、漏报、集合宽度和弃权成本写成显式参数，使不同风险偏好可以在同一框架下表达。

第六，审计优先。系统把证据、锚点、资格、可靠性、信念轨迹、动作损失和日期估计写入结构化文件。审计不是事后附加的日志，而是系统输出的一部分。没有审计文件的判断不能被本文视为可复核的判断。

\section{三层锚点状态}

FinPlanRAG的核心状态是AnchorEvolve的三层锚点。三层分别是稳定关系层（L0）、行业模式层（L1）和实例层（L2）。这里的锚点不是知识图谱中的实体或关系，也不是行为金融中的anchoring bias；L0、L1和L2分别对应关系约束、行业模式绑定和带来源的证据实例，锚定的是证据对结果的约束关系，而不是一个实体节点。AnE用“anchor”指训练数据的真值锚定，本文的锚点指证据状态中的结构化对象，二者机制不同\citep{wang2026ane}。

L0是稳定关系本体。它定义一组可复用的关系类型，包括界、先后、前置时间、持续时间、守恒、比率、单位换算和作用域。L0不随行业变化，为锚点提供语义骨架。例如，“当前季度值落在指引区间内”是一个界关系，“新产品在指引季度之后投产”是一个先后关系，“产能爬坡需要两个季度”是一个前置时间关系，“分部收入之和等于总收入”是一个守恒关系，“毛利率等于毛利润除以营业收入”是一个比率关系。L0的作用是限制锚点可以表达的关系种类，使不同行业的锚点共享同一套语义。

L1是行业模式。每个锚点类型是一个四元组
\begin{equation}
\tau=(\text{名称},\ \text{L0关系},\ \text{字段},\ \text{生效轮次}),
\label{eq:l1}
\end{equation}
其中字段给出该类型需要从证据中读取的信息，生效轮次给出该类型可以在哪一轮断言中使用。收入指引的初始模式包括指引宽度、当前季度增长、公司历史、行业历史、文本信号、产能与需求、分析师预测、披露历史和监管规则。L1把L0关系绑定到具体字段和生效规则。例如，“指引宽度”绑定到界关系，字段是上限与下限之差除以中点；“当前季度增长”绑定到比率关系，字段是当前季度值相对指引中点的比例；“披露历史”绑定到先后关系，字段是公司过去若干年年报的公告日期。

L2是锚点实例。每条锚点是一个元组
\begin{equation}
a=(id,\tau,v,u,s,src,d,r,\text{hard},\hat h,\text{note}),
\label{eq:l2}
\end{equation}
其中$id$是锚点编号，$\tau$是L1类型，$v$是取值或区间，$u$是单位，$s$是作用域，$src$是来源编号，$d$是来源日期，$r$是可靠性，$\text{hard}$是硬锚点标记，$\hat h$是隐含结果标签，$\text{note}$是理由。L2是LLM抽取器的直接输出，也是资格判定和信念更新的输入。每条L2锚点都记录来源和时间，因此可以被逐条检查。

三层的关系可以概括为：L0定义可以说哪些关系，L1定义某类断言需要哪些字段，L2记录某条证据的具体取值。三层不是三个独立训练的模块，而是同一状态的不同抽象层次。L1类型用于可靠性演化，L0关系用于语义分类，L2实例进入信念更新。这样的分层使系统可以在不改变L0的情况下增加新的L1类型，也可以在不改变L1的情况下替换L2抽取器。

\section{锚点类型}

表\ref{tab:anchortypes}给出本文使用的主要锚点类型。类型的选择遵循两个标准：一是该类型在金融文本中有稳定的语言表达，二是该类型能够对应到结果空间中的一个或多个标签。类型不是越多越好；过多的类型会增加抽取器的负担，也会增加可靠性后验的稀疏性。

\begin{table}[H]
\centering
\zihao{5}
\caption{FinPlanRAG的主要锚点类型。}
\label{tab:anchortypes}
\begin{tabular}{p{2.6cm}p{2.4cm}p{7.2cm}}
\toprule
类型 & L0关系 & 作用 \\
\midrule
指引宽度 & 界 & 指引区间越宽，结果落在区间内的先验越高；区间本身不是结果落在区间内的证据。\\
当前季度增长 & 比率 & 当前季度实际值相对指引中点的比例，用于判断全年结果的起点。\\
公司历史 & 先后 & 公司过去已揭示结果相对指引的偏差，用于估计公司层面的兑现倾向。\\
行业历史 & 先后 & 同行业其他公司已揭示结果的分布，用于估计行业层面的共同冲击。\\
文本信号 & 语义 & 管理层评论、需求描述、订单和产能信息，用于补充结构化字段。\\
产能与需求 & 前置时间 & 产能投放、利用率和订单覆盖，用于判断收入增长的可持续性。\\
分析师预测 & 界 & 卖方机构对目标期间的预测，作为软锚点而不是事实。\\
披露历史 & 先后 & 公司过去年度报告的公告日期，用于披露日期断言。\\
监管规则 & 界 & 交易所的披露截止日，作为上界而不是具体日期证据。\\
\bottomrule
\end{tabular}
\end{table}

锚点类型与结果标签的关系不是一一对应。一条“当前季度增长”锚点可能支持“高于”标签，也可能只支持“不低于”某个下界；一条“指引宽度”锚点通常不直接支持任何标签，而是影响区间内的先验。资格判定和信念更新需要分别处理这些情况。锚点的$\hat h$字段记录抽取器认为该锚点隐含的标签；如果锚点不提供标签信息，$\hat h$为空，锚点被判为空泛。

\section{证据来源与作用域}

FinPlanRAG使用三类证据来源。

第一类是公司申报文件和公告。美国公司使用8-K、10-Q、10-K和XBRL数据；中国A股公司使用巨潮资讯网的公告记录、年报、季报和临时公告。申报文件提供结构化事实和原文段落，公告记录提供披露日期和文件链接。官方来源的优先级高于媒体转载，因为官方来源的日期和文本更可靠。

第二类是财经媒体和研究报告。媒体文章可以提供管理层评论、行业需求、产能和价格信息，但可能存在转载、改编和日期错误。系统要求媒体来源给出逐字引文，并对无日期来源降权。研究报告中分析师预测作为软锚点，不作为事实。

第三类是交易所规则和监管文件。这类来源提供披露截止日、披露格式和合规要求。监管规则通常只给出上界，不能直接支持某个具体日期窗口；系统把泛泛的监管截止日判为空泛锚点，除非来源明确给出公司的预约披露日期。

锚点的作用域分为公司、行业、市场和监管四类。作用域的作用是防止跨作用域误用。例如，同行业其他公司的营收指引结果可以作为行业锚点，但不能直接当作目标公司的结果；交易所的披露截止日可以作为监管锚点，但不能直接推断公司会在截止日当天披露。资格判定检查作用域是否与断言匹配，不匹配的锚点被判为对抗锚点。

\section{与既有框架的区别}

FinPlanRAG与断言验证、RAG和选择性预测的区别可以概括为三点。

第一，状态不同。断言验证的输出是一个离散裁决，RAG的输出是一个生成答案，选择性预测的输出是一个标签或弃权。FinPlanRAG的输出是一个可演化的锚点状态，包含证据、资格、可靠性和信念。这个状态可以被逐条检查，也可以在结果揭示后更新。

第二，时间不同。断言验证和RAG通常在一次调用中完成，不区分证据的可用时间和结果的揭示时间。FinPlanRAG把截止时点写进契约，把揭示日写进可靠性更新规则。这个区分使系统可以在回测中复现部署时的信息条件。

第三，决策不同。断言验证和RAG通常报告准确率或F1，选择性预测报告覆盖率或弃权率。FinPlanRAG在同一个损失下比较点答案、裁决集和弃权，报告决策损失、覆盖率和集合大小。这个口径使系统可以在不确定性高时选择集合或弃权，而不是被迫输出一个点答案。

% ==================== 第5章 ====================

本章给出FinPlanRAG中AnchorEvolve模块的完整算法设计。本章第1节到第7节分别介绍证据构造、锚点抽取、资格判定、时间演化、信念更新、全动作读出和披露日期断言，最后给出算法伪代码、复杂度和审计格式。

\section{证据构造}

证据构造的目标是为每条断言提供一组带日期、带来源编号、可引用的材料。构造过程分为查询生成、检索、去重、日期过滤和编号五步。

查询生成从契约出发。给定公司$e$、目标期间$p$和指标$m$，系统生成四类查询：公司加期间加指标，公司加最近季度加指标，公司加上一年度加指标，以及公司加“预测”“研报”“公告”等限定词。对于披露日期断言，系统额外生成公司加“年度报告披露日期”、公司加“往年披露时间”、公司加“预约披露”等查询。查询数量由配置限制，默认每条断言最多四条。查询中不放入断言的目标答案，避免把判断泄漏到检索。

检索使用多个网页搜索提供者，并按轮换顺序调用。某个提供者返回空结果或安全验证页时，系统自动切换到下一个提供者。这个设计是为了避免单一搜索引擎的限流让整次运行静默失败。检索结果包含标题、链接、摘要、发布日期和提供者。对于披露日期断言，系统还查询巨潮资讯网的公告接口，直接获取公司历年年报的公告日期和PDF链接。官方公告记录作为高优先级证据，排在网页结果之前。

去重按链接和标题进行。同一篇文章的不同转载只保留一条；同一链接的不同摘要只保留一条。日期过滤检查证据的发布日期是否晚于截止时点。晚于截止时点的证据被写入剔除列表，不进入锚点抽取。发布日期未知的证据保留，但标记为无日期，并在资格判定中降权。

编号使用$W1,W2,\dots$，按证据得分排序。得分由日期是否已知、摘要是否包含数字、标题或摘要是否包含公司名、文本长度和查询顺序决定。官方公告记录得分最高。编号的作用是让锚点可以引用具体来源，也让审计报告可以把锚点映射回证据。

对于申报文件，系统在证据构造阶段还做段落切分。每个段落带稳定编号$P1,P2,\dots$，结构化证据带$E1,E2,\dots$。抽取器可以引用$P\#$或$E\#$，资格判定检查引用的段落是否在时点摘录中。这个设计使锚点可以定位到段落，而不是只定位到文件。

\section{锚点抽取}

锚点抽取是一个LLM调用。输入包括断言契约、结果标签及其定义、参考口径、时点证据、锚点类型和关系类型。输出是一个JSON对象，包含锚点列表、标签分数、日期估计（仅日期断言）、置信度和缺失说明。

提示词要求锚点遵守以下规则。第一，每条锚点必须引用一个证据编号，不能引用不存在的来源。第二，每条锚点必须给出逐字引文，引文必须在来源文本中出现。第三，锚点的发布日期不得晚于截止时点；来源日期未知时，日期设为空，可靠性不得高于0.55。第四，可靠性是“隐含标签正确的校准概率”，不是推理置信度。第五，硬锚点必须是显式约束，能够排除至少一个标签；管理层评论和分析师预测通常是软锚点。第六，证据不足时在缺失字段中说明，不编造事实。第七，标签分数必须使用契约中的标签，且和为1。

抽取器的输出是一个结构化对象，而不是自由文本。这个约束使后续的资格判定可以按字段检查，而不是依赖另一个LLM来判断前一个LLM的输出是否合理。解析失败的响应被记录为错误，不进入信念更新；系统不使用金标结果作为回退。这个设计是为了避免在评估中把解析失败悄悄地替换成正确答案。

锚点抽取的成本主要来自输入和输出长度。单条断言的输入通常在数千token量级，输出在数百token量级。系统按请求负载的哈希缓存响应，相同输入产生相同输出。缓存不包含API密钥，也不写入审计文件的公开部分。

\section{五类资格判定}

资格判定是FinPlanRAG与普通RAG的关键区别。普通RAG把检索到的片段直接交给生成器；FinPlanRAG在锚点进入信念更新之前，先用确定性规则检查锚点的资格。表\ref{tab:qualification}给出五类判定。

\begin{table}[H]
\centering
\zihao{5}
\caption{五类锚点资格判定。}
\label{tab:qualification}
\begin{tabular}{p{2.1cm}p{7.0cm}p{4.2cm}}
\toprule
类型 & 判定条件 & 处理 \\
\midrule
空泛 & 没有隐含标签，或锚点不提供有效标签信息 & 不进入更新 \\
过期 & 锚点日期或来源日期晚于截止时点 & 不进入更新 \\
对抗 & 来源编号不在证据集合中，或作用域不匹配 & 不进入更新 \\
硬 & 显式约束、引文通过、来源日期已知、可靠性不低于阈值 & 约束可行集 \\
软 & 其余有效锚点 & 似然更新 \\
\bottomrule
\end{tabular}
\end{table}

资格判定的具体规则如下。

第一，来源检查。锚点的来源编号必须出现在本次证据集合中。来源不存在，判为对抗。这个检查防止LLM引用一个看起来合理但实际不存在的来源。

第二，日期检查。锚点日期或来源日期晚于截止时点，判为过期。如果锚点日期与来源日期不一致，系统以来源日期为准检查过期，并对可靠性乘以0.9。日期不一致本身不是拒绝理由，但它是一个时间一致性疑点。

第三，引文检查。锚点的引文必须在来源文本中逐字出现。系统对引文和来源文本做空白规范化，然后做子串匹配。引文为空、长度过短或匹配失败，判为引文未通过。未通过的锚点可靠性乘以0.7，且不能成为硬锚点。这个检查是确定性的，不依赖另一个LLM。

第四，作用域检查。锚点的作用域必须是公司、行业、市场或监管之一。作用域不匹配，判为对抗。例如，把同行业其他公司的结果当作目标公司的直接结果，作用域不匹配。

第五，无日期检查。锚点日期和来源日期都未知时，锚点保留为软锚点，但可靠性上限为0.55，且不能成为硬锚点。这个规则承认无日期证据可能有用，但不允许它主导判断。

第六，参考值检查。对于增速类断言，只描述上一年度参考值的锚点判为空泛。上一年度实际值是契约的参考口径，不是目标期间的结果证据。这个检查防止抽取器把基数当作结果，从而把信念拉向低增长。

第七，监管规则检查。对于披露日期断言，泛泛的法定截止日判为空泛。截止日是一个上界，不能区分三个日期窗口；只有公司预约披露日期或历史披露记录才能支持具体窗口。

第八，日期标签投影。对于披露日期断言，如果锚点提到的日期与隐含标签不一致，系统把日期中的月和日投影到目标年份，再匹配日期窗口。例如，历史年报日期2025-04-21投影到2027-04-21，落在“晚于2027-04-15”窗口。投影只在一个锚点提到的所有日期指向同一窗口时生效；如果锚点同时提到多个窗口的日期，系统不覆盖模型标签，只记录不一致。这个规则修正了抽取器在跨年日期上的常见错误。

第九，时间一致性检查。对于披露日期断言，如果引文中提到的日期晚于来源发布日期，且不晚于截止时点，系统标记时间一致性疑点，可靠性乘以0.5，且不能成为硬锚点。这个检查处理动态网页或转载文章在原始发布日期之后被更新，导致引文包含当时不可能知道的信息的情况。

资格判定的输出是每条锚点的状态、资格类型、有效可靠性、引文是否通过、时间一致性标记和理由。被拒绝的锚点保留在审计文件中，理由逐条列出；被采纳的锚点进入信念更新。这个设计使审计者可以看到系统拒绝了什么，以及为什么拒绝。

\section{时间演化}

时间演化的目标是让锚点类型的可靠性随已揭示结果更新，同时不把未来信息泄漏到时点判断。设断言$q$的截止时点为$c$，锚点类型为$\tau$。截至$c$已揭示的、类型为$\tau$的锚点中，隐含标签与结果一致的数量为$m_\tau$，总数为$n_\tau$。类型$\tau$的Beta可靠性为
\begin{equation}
r_\tau=\frac{m_\tau+\alpha}{n_\tau+\alpha+\beta},
\label{eq:beta}
\end{equation}
其中$\alpha,\beta$是Beta先验参数。默认取$\alpha=\beta=1$，即均匀先验。当$n_\tau=0$时，$r_\tau=0.5$。

抽取器给出的可靠性是LLM对单条锚点的主观判断，不是经过验证的频率。系统把演化可靠性与LLM可靠性做收缩：
\begin{equation}
\rho_\tau=\lambda r_\tau+(1-\lambda)r_{\mathrm{LLM}},
\label{eq:shrink}
\end{equation}
其中$\lambda$是收缩系数，默认0.5。$\rho_\tau$作为软锚点更新的指数。这个收缩的作用是：当历史样本少时，系统更多依赖LLM判断；当历史样本多时，系统更多依赖已揭示结果的频率。

先验也按揭示日更新。设截至$c$已揭示的结果中，落在标签$h$的数量为$n_h$，总数为$n$，则先验为
\begin{equation}
b_0(h)=\frac{n_h+1}{n+|\mathcal H|}.
\label{eq:prior}
\end{equation}
当$n=0$时，先验为均匀分布。这个先验只使用已揭示结果，不使用目标期间的结果。

演化状态按轮次更新。第$r$轮的结果在第$r$轮全部断言裁决完成后才更新状态，并从第$r+1$轮开始使用。新发现的锚点类型的生效轮次为$r+1$，因此触发发现的断言不能使用该类型。这个规则保证类型库扩展是时点正确的。本文的主实验报告单轮扩展；多轮扩展作为接口保留，第\ref{ch:conclusion}章讨论其局限。

\section{类型库的验证集扩展}

固定锚点类型库依赖人工设计，难以覆盖不同指标和不同市场的证据模式。本文在Beta可靠度演化之外加入一个受验证集约束的类型库扩展步骤。它不是让模型自由改写标签或测试集，而是让模型在验证误差上提出新的锚点类型，只有通过验证集损失检验的类型才进入类型库。

设第$t$轮的锚点类型库为$\mathcal T_t$，验证集上由当前类型库得到的决策损失为$L_{\mathrm{val}}(\mathcal T_t)$。循环分四步。

第一步，误差挖掘。把验证集按公司不重叠地分成挖掘集和选择集。在挖掘集上，用当前类型库计算每条断言的预测和损失，选出预测错误或损失较高的断言。挖掘集只用于提出候选类型，不参与候选接受。

第二步，类型提出。把错误断言、对应证据、当前锚点和真实结果交给一个LLM，要求它提出至多$k$个新的锚点类型。每个候选类型是一个四元组
\begin{equation}
\tau'=(\text{名称},\ \text{L0关系},\ \text{字段},\ \text{适用条件}),
\label{eq:candidatetype}
\end{equation}
并附带一句证据示例。提出器被要求只描述可复用的证据模式，不能引用具体测试断言，不能把结果标签本身当作类型。

第三步，候选重抽取。把候选类型加入类型库，重写锚点抽取提示词中的类型说明，在选择集上重新抽取锚点并计算后验。类型库扩展为$\mathcal T_t\cup\mathcal T'$。

第四步，验证门控。只有当选择集上的损失下降超过阈值$\delta$时，才接受候选类型：
\begin{equation}
L_{\mathrm{val}}(\mathcal T_t)-L_{\mathrm{val}}(\mathcal T_t\cup\mathcal T')>\delta.
\label{eq:typeaccept}
\end{equation}
否则拒绝整组候选，保持原类型库。接受后进入下一轮，最多进行$R$轮。测试集只在循环冻结后评估一次，不参与类型提出、阈值选择或轮数选择。

这个循环与提示词演化分开评估。提示词改写、类型扩展和二者组合分别作为候选臂，在同一选择集上比较；只有损失最低且超过阈值的臂被接受。这样可以把“模型提出了新类型”和“新类型真的改善了验证损失”分开，避免把语言上的合理性当作经验收益。

类型库扩展的代价是额外的验证集抽取。每轮需要重抽选择集锚点，但不重抽测试集；类型库冻结后才抽取测试集。由于循环只使用验证集，它与主实验的时点约束一致，也不会把测试结果泄漏到类型库。

类型库扩展是框架的训练步骤，不是独立模块。算法\ref{alg:typeexpand}给出伪代码，其输出$\mathcal T^\star$作为算法\ref{alg:anchorevolve}的类型库输入。主实验中DeepSeek使用通过验证门控的类型库，GLM和Qwen保留基础库；这样，类型库的选择和其他超参数一样，只由验证集决定。

\begin{algorithm}[H]
\caption{类型库的验证集扩展}
\label{alg:typeexpand}
\begin{algorithmic}[1]
\Require 类型库$\mathcal T_0$，验证集$\mathcal D_{\mathrm{val}}$，阈值$\delta$
\Ensure 冻结类型库$\mathcal T^\star$
\State 将$\mathcal D_{\mathrm{val}}$按公司不重叠划分为挖掘集$\mathcal D_{\mathrm{mine}}$和选择集$\mathcal D_{\mathrm{select}}$
\State 在$\mathcal D_{\mathrm{mine}}$上计算当前类型库的损失，选出高损失断言
\State 由LLM提出候选类型$\mathcal T'$，每个类型包含关系、字段和适用条件
\State 在$\mathcal D_{\mathrm{select}}$上重抽锚点，计算$L_{\mathrm{val}}(\mathcal T_0\cup\mathcal T')$
\If{$L_{\mathrm{val}}(\mathcal T_0)-L_{\mathrm{val}}(\mathcal T_0\cup\mathcal T')>\delta$}
\State $\mathcal T^\star\gets\mathcal T_0\cup\mathcal T'$
\Else
\State $\mathcal T^\star\gets\mathcal T_0$
\EndIf
\State \Return $\mathcal T^\star$ \Comment{冻结后用于测试集}
\end{algorithmic}
\end{algorithm}


\section{信念更新}

信念从先验$b_0$开始。对于软锚点$a$，设其隐含标签为$\hat h$，有效可靠性为$\rho$，软似然为
\begin{equation}
L_a(h)=
\begin{cases}
p, & h=\hat h,\\
\dfrac{1-p}{|\mathcal H|-1}, & h\ne\hat h,
\end{cases}
\label{eq:softlikelihood}
\end{equation}
其中$p$是配置项，默认0.7。软锚点的更新为
\begin{equation}
b(h)\propto b(h)\,L_a(h)^{\rho}.
\label{eq:softupdate}
\end{equation}
指数$\rho$控制锚点的影响强度：可靠性越高，更新越强；可靠性越低，更新越弱。更新后对信念做归一化。

对于硬锚点，系统把信念压到隐含标签上，并保留一个小底：
\begin{equation}
b(h)=
\begin{cases}
1-\epsilon, & h=\hat h,\\
\dfrac{\epsilon}{|\mathcal H|-1}, & h\ne\hat h,
\end{cases}
\label{eq:hardupdate}
\end{equation}
其中$\epsilon$默认0.02。小底的作用是避免后续软锚点无法恢复。多个硬锚点的可行集交集、冲突检测和回退到软更新是论文的设计目标，当前实现只处理单硬锚点的点质量更新；第\ref{ch:conclusion}章把这一点列为局限。

锚点后验记为$b_A$。抽取器还输出标签分数$s$，它是边际信号，不是锚点状态的替代。系统把$b_A$与$s$按权重$\gamma$混合：
\begin{equation}
b_A'=\gamma b_A+(1-\gamma)s,
\label{eq:blend}
\end{equation}
其中$\gamma$默认0.5。混合后的后验保留锚点的结构化信息，同时利用抽取器对三类结果的边际判断。

结构化随机森林后验记为$b_R$。输入特征包括指引宽度、当前季度值相对指引中点的比例、指引与当前季度的比值、历史结果计数与频率、前次指引增长与宽度、目标季度、截止年份、行业，以及演化先验特征（同行业历史、公司意外历史、日历、当前季度相对上下界的位置、可靠性标志）。随机森林在开发集与验证集上训练，并在公司不重叠折上做温度缩放；测试公司不出现在训练集中。

锚点后验与随机森林后验融合：
\begin{equation}
b_F=w b_A'+(1-w)b_R,
\label{eq:fusion}
\end{equation}
其中$w$只在验证集上从配置的权重网格中选择，测试集不参与选择。默认网格为$\{0,0.1,0.2,0.3\}$。最终实验中选择的权重在正文中报告。

\section{全动作损失感知读出}

给定融合后验$b_F$，系统在动作空间$\mathcal A$上选择期望损失最小的动作。动作空间包括$|\mathcal H|$个点答案、$\sum_{k=2}^{|\mathcal H|}\binom{|\mathcal H|}{k}$个非单元素裁决集，以及弃权。对于三标签结果空间，动作空间包含3个点答案、4个非单元素裁决集和1个弃权，共8个动作。

点答案的期望损失为
\begin{equation}
\E[L_{\text{ans}}(h)]=\sum_{y\in\mathcal H}b_F(y)L_{\text{ans}}(h,y).
\label{eq:eans}
\end{equation}
裁决集的期望损失为
\begin{equation}
\E[L_{\text{set}}(R)]=\sum_{y\in\mathcal H}b_F(y)L_{\text{set}}(R,y).
\label{eq:eset}
\end{equation}
弃权的期望损失为$\ell_{\text{abstain}}$。系统选择
\begin{equation}
a^\star=\arg\min_{a\in\mathcal A}\E[L(a)].
\label{eq:fullaction}
\end{equation}

全动作读出与受限可信集门的区别在于动作空间。受限门只比较点答案和最高密度$1-\alpha$可信集，可信集的大小由覆盖率目标决定，而不是由损失决定。全动作读出把点答案、裁决集和弃权放在同一个损失下比较，因此可以在不确定性高时选择裁决集或弃权，也可以在置信度高时选择点答案。对于同一个后验和同一个损失，全动作读出的期望损失不会高于受限门，因为受限门的动作是全动作空间的子集。这个性质是H3的形式基础。

单元素裁决集不进入动作空间。原因是单元素裁决集与点答案在语义上重复，但在损失上不同：如果漏报惩罚低于错误分类惩罚，单元素裁决集会系统性地优于点答案，导致系统把点答案都标成裁决集。这个现象不是决策改进，而是动作标签的重复计数。本文把裁决集限制为两个及以上标签，使三类动作的含义保持清晰。

可信集$R_\alpha$仍然作为覆盖口径报告：
\begin{equation}
R_\alpha=\arg\min_{R\subseteq\mathcal H}\{|R|:\sum_{h\in R}b_F(h)\ge1-\alpha\}.
\label{eq:credible}
\end{equation}
它不参与动作选择，但用于审计和覆盖率报告。损失最优裁决集$R^\star$是动作选择的结果，可以在$R_\alpha$更宽时选择更窄的集合，也可以在不确定性高时选择弃权。两者在审计报告中分开列出，避免把覆盖率误读为决策质量。

\section{披露日期断言}

披露日期断言的算法与数值指标断言共享同一套锚点状态和读出，但结果空间和日期估计不同。

结果空间由日历窗口定义。给定目标期间$p$，系统根据交易所规则和公司历史披露习惯构造三个窗口。例如，2026年年度报告的窗口为“不晚于2027-03-31”“2027-04-01至2027-04-15”“晚于2027-04-15”。窗口必须互斥且覆盖法定披露截止日之前的全部日期。

锚点抽取器除了输出锚点和标签分数，还输出一个日期估计，包括点估计、区间下界、区间上界和依据。日期估计是辅助信息，不直接决定动作；动作仍由标签后验和损失读出决定。系统把日期估计与标签后验做一致性检查：如果点估计落在最高信念窗口内，系统保留点估计；如果点估计落在窗口外，系统用最高信念窗口的中点作为回退点估计。区间下界和上界优先使用模型给出的值，但限制在合理范围内。

日期标签投影是披露日期断言的关键修正。历史年报日期、预约披露日期和监管截止日往往跨年，抽取器容易把历史日期直接映射到当前窗口。系统把日期中的月和日投影到目标年份，再匹配窗口。投影规则只在一个锚点提到的所有日期指向同一窗口时生效，避免把包含多个日期的锚点强行归类。

披露日期断言的评价不以点日期误差作为主指标，因为公告日期与实际披露时点可能相差一天。本文在A股回测上报告日期窗口准确率，并在案例研究中展示动作、裁决集、点估计、区间和审计轨迹。

\section{算法与复杂度}

算法\ref{alg:anchorevolve}给出单条断言的完整流程。

\begin{algorithm}[H]
\caption{AnchorEvolve单条断言的证据构造、锚点演化与全动作裁决}
\label{alg:anchorevolve}
\begin{algorithmic}[1]
\Require 断言契约$q=(e,p,m,\mathcal H,c)$，演化状态$S_c$，类型库$\mathcal T^\star$，随机森林后验$b_R$
\Ensure 动作$a^\star$，点答案$h^\star$，裁决集$R^\star$，后验$b_F$
\State $V_c\gets\textsc{CollectEvidence}(q,c)$ \Comment{按截止时点过滤，分配$W\#$}
\State $A\gets\textsc{ExtractAnchors}(q,V_c,S_c,\mathcal T^\star)$ \Comment{LLM，要求来源与逐字引文}
\State $A_q\gets\textsc{Qualify}(A,q,c)$ \Comment{空泛/过期/对抗/硬/软}
\State $b_A\gets\textsc{BeliefUpdate}(S_c,A_q)$
\State $b_A'\gets\gamma b_A+(1-\gamma)\textsc{LabelScores}(A)$
\State $b_F\gets w b_A'+(1-w)b_R$
\State $(a^\star,h^\star,R^\star)\gets\textsc{FullActionReadout}(b_F)$
\State $S_{c^+}\gets\textsc{RevealUpdate}(S_c,c)$ \Comment{仅用揭示日早于下一截止时点的结果}
\State \Return $(a^\star,h^\star,R^\star,b_F)$
\end{algorithmic}
\end{algorithm}

算法的时间成本由三部分组成。第一是证据构造，包括网页检索、页面抓取和官方公告查询。网页检索的数量由配置限制，默认每条断言四条查询；页面抓取只抓前若干条结果，避免不必要的网络开销。第二是锚点抽取，一次LLM调用的输入在数千token量级，输出在数百token量级。第三是读出，动作空间的大小为$|\mathcal H|+\sum_{k=2}^{|\mathcal H|}\binom{|\mathcal H|}{k}+1$，当$|\mathcal H|\le6$时最多63个动作，计算可以忽略。随机森林推断在本地完成，不产生API成本。

系统的存储成本主要来自审计文件。每条断言保存证据列表、锚点列表、信念轨迹、动作损失和LLM响应摘要。审计文件的大小通常在数百KB量级，适合长期保存和人工复核。LLM响应缓存单独存放，不进入公开审计文件。

\section{可审计性}

可审计性是FinPlanRAG的设计目标之一。系统把每个中间状态写入结构化文件，使读者可以从最终动作回溯到证据和锚点。表\ref{tab:auditfields}给出审计文件的主要字段。

\begin{table}[H]
\centering
\zihao{5}
\caption{审计文件的主要字段。}
\label{tab:auditfields}
\begin{tabular}{p{2.8cm}p{8.8cm}}
\toprule
对象 & 字段 \\
\midrule
契约 & 公司、代码、目标期间、指标、结果空间、标签定义、参考口径、截止时点 \\
证据 & 来源编号、标题、链接、发布日期、提供者、摘要、正文、是否官方来源 \\
锚点 & 类型、取值、来源编号、逐字引文、发布日期、可靠性、资格、隐含标签、理由、纠正记录 \\
信念 & 先验、仅锚点后验、融合后验、逐步更新轨迹 \\
动作 & 动作类型、点答案、裁决集、可信集、期望损失、各动作损失排序 \\
日期 & 点估计、区间、最高信念窗口、模型估计与窗口是否一致、依据 \\
\bottomrule
\end{tabular}
\end{table}

审计文件有两个用途。第一，复核。分析师可以逐条检查被采纳的锚点是否真的支持隐含标签，被拒绝的锚点是否应当被拒绝，证据的发布日期是否早于截止时点。第二，问责。当判断出现错误时，系统可以定位错误来自证据缺失、锚点抽取、资格判定、信念更新还是动作选择，而不是把错误归因于一个黑箱模型。

可审计性不等于正确性。一个锚点即使有来源和引文，也可能是对来源的误读；一个资格判定即使规则明确，也可能过严或过松。本文在案例研究中展示审计轨迹，并在第\ref{ch:conclusion}章把人工审计实验列为未来工作。当前的可审计性是设计属性，不是已经测量的用户收益。


% ==================== 第6章 ====================
\chapter{实证结果与分析}
\label{ch:empirical}

\section{数据与任务}

本文在四个数据集和一个混合负载上评估FinPlanRAG。四个数据集分别对应长文档确定问答、前瞻性断言裁决、外部迁移和A股年报回测，混合负载用于检验路由。这样的设计使框架的不同主张分别接受检验，而不是用一个数据集的平均结果替代全部证据。

DocFinQA在FinQA问题的基础上加入完整10-K上下文\citep{reddy2024docfinqa}。本文使用测试集中的300题固定样本，答案使用发布的金标准。FinQA程序常量只作为证据召回率的银标准，不作为阅读器输入。该数据集的挑战在于证据分散在长文档中，表格行与段落需要联合阅读，阅读器必须在大量无关内容中定位分子、分母和期间。

收入指引基准包含2010年至2025年美国SEC 8-K文件中的1,005条季度营收指引断言。每条断言包含指引区间、截止时点、当前季度值和证据项。实际季度营收来自SEC XBRL，XBRL没有独立季度时使用披露文件回退。标签为低于区间、区间内、高于区间。数据按时间划分为529条开发集（2010年至2018年）、203条验证集（2019年至2021年）和273条测试集（2022年至2025年）。测试集包含126条高于、110条区间内和37条低于。主结果使用227条具有XBRL核验或XBRL覆盖标签的测试断言；46条纯文本标签单独用于敏感性检验。该数据集的挑战在于标签来自自动核验而非人工标注，类别不平衡明显，且前瞻性断言的真值在截止时点尚未实现。

FinDVer testmini包含700条带财务报告的断言验证样本\citep{zhao2024findver}，作为外部迁移集。本文不在FinDVer上调参，只把在DocFinQA上设计的义务检索策略迁移过去，以检验证据策略是否依赖特定数据分布。

第四个基准是A股年报历史回测。本文选取40家非金融上市公司，覆盖消费、制造、医药、能源和科技行业，对2023年至2025年三个财年构造470条断言，其中营业收入增速232条、归母净利润增速120条、综合毛利率120条、披露日期118条。对每个财年$Y$，问题在$Y+1$年1月31日提出，此时三季报已经披露、年报尚未披露；结果来自随后披露的年报。该基准检验模糊分支在中国市场上的时点判断、数值桥和披露日期断言。

混合负载由300条DocFinQA问题和273条收入指引断言组成，共573条查询。确定查询的效用为词级F1；前瞻性查询的效用为$1-L/10$，其中$L$为决策损失。错误路由记0效用。Oracle路由器把每条查询送到正确分支；Always-RAG把所有查询送到义务检索分支；Always-Anchor把所有查询送到AnchorEvolve。混合效用把两类任务的分数放在同一尺度上，使路由策略可以与固定分支策略比较。

\section{基座模型、具名基线与评价指标}

本文使用三套基座：DeepSeek Flash、GLM 5.3 Flash和Qwen3-8B\citep{qwen2025qwen3}。Qwen3-8B在本地部署。所有方法使用同一基座和相同的提示模板，以比较证据策略和输出契约的差异。

义务检索分支与固定BM25预算1、2、4、8、16段、Self-RAG、CRAG、FLARE、Adaptive-RAG以及义务检索规划器比较。四个具名RAG方法根据原论文描述和公开接口进行同基座适配，保留各自的检索或反思策略，使用与义务检索分支相同的段落块、阅读器和评分，因此差异来自证据策略。

AnchorEvolve与六个具名方法比较：FinDVer \citep{zhao2024findver}、Fin-Fact \citep{rangapur2025finfact}、Self-RAG \citep{asai2024selfrag}、CRAG \citep{yan2024crag}、FLARE \citep{jiang2023flare}和Adaptive-RAG \citep{jeong2024adaptiverag}。每个方法根据原论文描述和公开接口进行同基座适配，通过同一个三分类适配层运行。适配层把原始输出映射为低于、区间内、高于，并提供相同的时点证据块和相同的六条历史断言。学习组件不可用或不兼容时，用使用官方判据的API调用替代，并记录替代关系。六个条件保留各自原论文的判据，不是原系统在其原始数据上的未修改端到端运行。

评价指标分为三组。DocFinQA报告精确匹配、词级F1、答案覆盖率、平均获取段落数、平均输入token和银标准证据召回率。收入指引报告准确率、平衡准确率、宏F1、MCC、AUROC、AUPRC、ECE、Brier、决策损失、可信集覆盖率和平均集合大小。A股回测报告准确率、平衡准确率、宏F1、决策损失、数值MAE、数值MAPE和日期窗口准确率。平衡准确率是各类召回率的平均\citep{brodersen2010balanced,sokolova2009performance}。决策损失使用第\ref{ch:anchor}章参数。损失、ECE和Brier越低越好。置信区间由断言级配对自助法给出，准确率差异使用McNemar检验。路由报告路由准确率、混淆矩阵和混合效用。

评价协议在三个层面控制信息泄漏。第一，收入指引按时间划分，测试年份为2022年至2025年，验证集为2019年至2021年，开发集为2010年至2018年。第二，随机森林使用五个折种子；对每个测试年份$Y$，训练集来自其他公司折，且结果日期早于$Y$年1月1日；测试公司不出现在训练集中。温度缩放在公司不重叠的内部折上拟合。第三，融合权重$w$只在验证集上按决策损失选择，候选为$\{0,0.05,0.1,0.2,0.3,0.5,1.0\}$；最终读出统一使用全动作门，不再在随机森林头与融合头之间切换。测试集不参与拟合、校准、权重选择或动作选择。

所有LLM响应按请求负载的哈希缓存。API密钥只从环境变量读取，不写入结果文件。提示词、运行脚本、逐条预测和表格生成脚本随论文发布。随机森林的折种子、校准方式和融合权重选择在正文中给出。公开结果包含逐条预测、路由决策和汇总指标；LLM响应缓存不随论文提交，但相同输入可以用相同提示词复现。

\section{路由与混合负载}

表\ref{tab:router}报告573条混合负载。DeepSeek Flash路由器正确路由570条（0.995）：300条确定问题中297条，273条前瞻性断言全部正确。GLM 5.3 Flash路由器正确路由567条（0.990）：297条确定问题、270条前瞻性断言。Qwen3-8B正确路由271条确定问题和226条前瞻性断言（0.867），其中29条确定问题被送到AnchorEvolve，47条前瞻性断言被送到义务检索，显示8B基座对契约边界的识别较不稳定。

DeepSeek Flash的混合效用为0.403，与oracle相同，高于Always-RAG的0.121和Always-Anchor的0.282。GLM 5.3 Flash的混合效用为0.419，接近oracle的0.422，并恢复99.3\%的oracle效用。Qwen3-8B的混合效用为0.263，低于Always-Anchor的0.298。结果表明，较大基座能够把确定问答与前瞻性裁决送入相应接口，小基座则需要更强的路由监督或外部信号。

路由错误主要出现在带有未来时态的确定问题，例如“未来支付金额”和“预期股利”。这些问题的答案已经在截止时点文件中披露，但表面表达接近前瞻性断言。该现象说明契约识别依赖事实是否已经存在，而不能仅依据时态或关键词。固定分支策略只能对全部查询采用同一接口，路由器则把契约选择显式化，使路由错误与分支错误可以分别诊断。

\begin{table}[H]
\centering
\zihao{5}
\setlength{\tabcolsep}{3pt}
\caption{混合负载路由结果。效用定义见第\ref{ch:empirical}章。}
\label{tab:router}
\begin{tabular}{lrrrrr}
\toprule
基座 & 路由准确率 & 路由器 & 始终RAG & 始终Anchor & Oracle \\
\midrule
DeepSeek Flash & .995 & .403 & .121 & .282 & .403 \\
GLM 5.3 Flash & .990 & .419 & .130 & .292 & .422 \\
Qwen3-8B & .867 & .263 & .018 & .298 & .316 \\
\bottomrule
\end{tabular}
\end{table}

\section{义务检索在DocFinQA上的结果}

表\ref{tab:docfinqa}报告DocFinQA的300题固定样本。最终义务检索分支使用16段预算，并在阅读器不充分或置信度低于0.8时回退到固定16段证据集。在DeepSeek Flash上，该分支取得0.457精确匹配，平均使用10.69段；固定16段阅读器取得0.390，平均16段；固定8段阅读器取得0.367，平均8段。无回退的预算8规划器取得0.417，平均4.56段；预算12取得0.427，平均5.82段；预算16取得0.460，平均6.85段。回退在24\%的问题上使用更大证据集，弥补了规划器单独的准确率缺口。

在GLM 5.3 Flash上，该分支取得0.497，平均10.31段；固定16段阅读器取得0.480，平均16段；固定8段阅读器取得0.440，平均8段。无回退的预算8规划器取得0.450，平均4.15段；预算12取得0.473，平均5.11段；预算16取得0.470，平均5.88段。回退分支比固定16段阅读器高0.017，平均少用5.69段。两个基座上的答案覆盖率分别为0.770和0.810，高于固定预算的0.640和0.757，说明分支更少弃权。

义务检索分支在准确率和证据成本上优于固定16段基线：准确率更高、证据更少。规划器单独运行时成本更低，但在GLM 5.3 Flash上留下小缺口；置信度门控回退在不对每个问题支付16段成本的情况下补上了缺口。这一结果的含义是证据预算应当由证据状态决定，而不是由查询前固定的$k$决定。对于简单问题，规划器在少量证据后停止；对于复杂问题，回退提供额外证据。

表\ref{tab:docfinqa}同时报告四个具名RAG适配基线。DeepSeek Flash上，Self-RAG与CRAG取得0.373，Adaptive-RAG取得0.363，FLARE取得0.330；GLM 5.3 Flash上，Self-RAG与CRAG取得0.447，Adaptive-RAG与FLARE取得0.393。义务检索分支在DeepSeek Flash上比最好的具名基线高0.084，在GLM 5.3 Flash上高0.050，且平均段落数不高于这些基线。四个条件保留各自的检索或反思策略，使用相同的段落块、阅读器和评分，因此差异来自证据策略而不是阅读器。

Qwen3-8B的DocFinQA结果放在附录\ref{app:qwen}：所有条件的精确匹配都低于0.11，主要原因是阅读器返回解释性文字而不是短答案，该基座不用于义务检索分支的主比较。该结果反映短答案抽取任务对阅读器格式遵循能力的要求，Qwen3-8B的阅读器更偏向解释性输出。框架的分支设计不依赖特定基座，但阅读器的格式遵循能力是实际部署的前提。

\begin{table}[htbp]
\centering
\zihao{5}
\setlength{\tabcolsep}{2pt}
\caption{DocFinQA长文档问答结果。Acc.为精确匹配，Macro F1为宏F1，Paras为平均段落，Cov.为覆盖率。}
\label{tab:docfinqa}
\begin{tabular}{lrrrrrrrr}
\toprule
 & \multicolumn{4}{c}{DeepSeek Flash} & \multicolumn{4}{c}{GLM 5.3 Flash} \\
\cmidrule(lr){2-5}\cmidrule(lr){6-9}
Condition & Acc. & Macro F1 & Paras & Cov. & Acc. & Macro F1 & Paras & Cov. \\
\midrule
Fixed-4 & .333 & .192 & 4.00 & .557 & .400 & .199 & 4.00 & .653 \\
Fixed-8 & .367 & .223 & 8.00 & .603 & .440 & .225 & 8.00 & .720 \\
Fixed-16 & .390 & .234 & 16.00 & .640 & .480 & .249 & 16.00 & .757 \\
Planner B=8 & .417 & .230 & 4.56 & .707 & .450 & .248 & 4.15 & .763 \\
Planner B=16 & .460 & .276 & 6.85 & .773 & .470 & .243 & 5.88 & .840 \\
Self-RAG & .373 & .223 & 8.03 & .613 & .447 & .231 & 8.16 & .717 \\
CRAG & .373 & .222 & 8.05 & .613 & .447 & .235 & 8.32 & .727 \\
FLARE & .330 & .197 & 4.00 & .557 & .393 & .193 & 4.19 & .643 \\
Adaptive-RAG & .363 & .211 & 7.22 & .570 & .393 & .198 & 6.51 & .657 \\
\midrule
\textbf{FinPlanRAG (obligation + fallback)} & \textbf{.457} & \textbf{.276} & 10.69 & \textbf{.770} & \textbf{.497} & \textbf{.252} & 10.31 & \textbf{.810} \\
\bottomrule
\end{tabular}
\end{table}


总体来看，义务检索分支在DeepSeek Flash和GLM 5.3 Flash上均以低于固定16段的平均证据量取得更高的准确率和覆盖率。规划器单独运行的成本最低，回退主要改善规划器无法闭合的难题；两者组合后的分支在准确率、覆盖率和证据成本之间取得较优的平衡。与四个具名RAG适配基线相比，义务检索的优势来自证据策略，而不是来自不同的阅读器或评分口径。

\section{前瞻性断言裁决结果}

表\ref{tab:main}按方法列出三套基座的结果。纵向是方法，横向是基座和指标，每个基座给出准确率、平衡准确率、宏F1和全动作决策损失。DeepSeek使用验证集接受的扩展类型库，GLM和Qwen保留基础类型库。

\begin{table}[htbp]
\centering
\zihao{-5}
\setlength{\tabcolsep}{2.2pt}
\caption{主结果（227条核验断言）。Acc.为准确率，Bal.为平衡准确率，Macro为宏F1，Loss为全动作决策损失。}
\label{tab:main}
\begin{tabular}{lrrrrrrrrrrrrr}
\toprule
\multirow{2}{*}{方法} & \multicolumn{4}{c}{DeepSeek Flash} & \multicolumn{4}{c}{GLM 5.3 Flash} & \multicolumn{4}{c}{Qwen3-8B} \\
\cmidrule(lr){2-5}\cmidrule(lr){6-9}\cmidrule(lr){10-13}
 & Acc. & Bal. & Macro & Loss & Acc. & Bal. & Macro & Loss & Acc. & Bal. & Macro & Loss \\
\midrule
FinDVer & .515 & .424 & .423 & 3.128 & .507 & .456 & .458 & 3.198 & .463 & .441 & .429 & 5.339 \\
Fin-Fact & .524 & .428 & .429 & 3.251 & .489 & .401 & .400 & 3.410 & .507 & .473 & .459 & 4.881 \\
Self-RAG & .546 & .460 & .468 & 3.031 & .507 & .454 & .460 & 3.181 & .485 & .446 & .431 & 5.119 \\
CRAG & .515 & .412 & .407 & 3.004 & .498 & .424 & .427 & 3.480 & .467 & .433 & .417 & 5.251 \\
FLARE & .511 & .418 & .412 & 3.278 & .489 & .407 & .408 & 3.559 & .326 & .298 & .212 & 6.661 \\
Adaptive-RAG & .489 & .402 & .396 & 3.612 & .436 & .366 & .366 & 4.070 & .352 & .325 & .290 & 6.396 \\
\midrule
LLM-only & .498 & .397 & .362 & 3.445 & .533 & .428 & .414 & 3.427 & .410 & .357 & .267 & 5.498 \\
RF & .533 & .491 & .486 & 2.987 & .533 & .491 & .486 & 2.987 & .533 & .491 & .486 & 2.987 \\
\textbf{FinPlanRAG (Ours)} & \textbf{.577} & \textbf{.526} & \textbf{.530} & \textbf{2.907} & \textbf{.551} & \textbf{.504} & \textbf{.503} & \textbf{2.943} & \textbf{.542} & \textbf{.501} & \textbf{.499} & \textbf{2.925} \\
\bottomrule
\end{tabular}
\end{table}

表中的FinPlanRAG指锚点后验与随机森林后验按验证集权重融合后的系统，融合权重只在验证集选择。DeepSeek的FinPlanRAG使用验证集扩展后的类型库，GLM和Qwen使用基础类型库；类型库扩展是训练步骤，不是另一个预测模块。FinPlanRAG在三个基座的准确率、平衡准确率、宏F1和决策损失上均为最高。DeepSeek的FinPlanRAG相对最强基线CRAG把损失从3.004降到2.907，准确率从.515提高到.577；GLM相对Self-RAG把损失从3.181降到2.943，准确率从.507提高到.551；Qwen相对Fin-Fact把损失从4.881降到2.925，准确率从.507提高到.542。配对自助法给出的损失差分别为$-0.097$ $[-0.502,0.300]$、$-0.238$ $[-0.661,0.176]$和$-1.956$ $[-2.617,-1.322]$，只有Qwen的区间排除零。

{\sloppy FinPlanRAG相对随机森林的损失差为$-0.079$ $[-0.300,0.141]$、$-0.044$ $[-0.203,0.115]$和$-0.062$ $[-0.282,0.159]$，三个区间均包含零。\par}融合在点估计上低于随机森林，但总体差异不足以主张稳定优势。把测试集按真实结果分层后，融合在区间内样本上的优势更清楚，见第\ref{sec:within}节。

表\ref{tab:calibration}给出概率质量指标。融合系统在三个基座上的AUROC为.713、.711和.707，ECE为.064、.030和.069，Brier为.555、.557和.561；对应的随机森林AUROC为.706，ECE为.033，Brier为.566。FinPlanRAG的排序能力和Brier分数优于随机森林；ECE在DeepSeek和Qwen上略高，在GLM上略低，说明融合后验的排序更准，但置信度校准仍有改进空间。

\begin{table}[htbp]
\centering
\zihao{5}
\setlength{\tabcolsep}{4pt}
\caption{概率质量与校准（227条核验断言）。AUROC和AUPRC越高越好，ECE和Brier越低越好。}
\label{tab:calibration}
\begin{tabular}{lrrrrrrrrr}
\toprule
\multirow{2}{*}{方法} & \multicolumn{3}{c}{DeepSeek Flash} & \multicolumn{3}{c}{GLM 5.3 Flash} & \multicolumn{3}{c}{Qwen3-8B} \\
\cmidrule(lr){2-4}\cmidrule(lr){5-7}\cmidrule(lr){8-10}
 & AUROC & ECE & Brier & AUROC & ECE & Brier & AUROC & ECE & Brier \\
\midrule
LLM-only & .668 & .110 & .612 & .651 & .101 & .609 & .493 & .439 & .975 \\
RF & .706 & .033 & .566 & .706 & .033 & .566 & .706 & .033 & .566 \\
\textbf{FinPlanRAG (Ours)} & \textbf{.713} & \textbf{.064} & \textbf{.555} & \textbf{.711} & \textbf{.030} & \textbf{.557} & \textbf{.707} & \textbf{.069} & \textbf{.561} \\
\bottomrule
\end{tabular}
\end{table}

\section{决策读出的作用}

表\ref{tab:readout}在同一后验上比较受限可信集门与全动作读出。受限门只比较点答案和最高密度可信集，全动作门在点答案、非单元素裁决集和弃权上选择期望损失最小的动作。融合后验在三个基座上的全动作损失相对受限门分别下降0.749、0.581和0.696，95\%配对区间均排除零。随机森林后验的下降为0.467，区间同样排除零。全动作读出没有改变后验，也没有增加训练数据，它改变的是给定后验时的动作选择。

\begin{table}[htbp]
\centering
\zihao{-5}
\setlength{\tabcolsep}{2.5pt}
\caption{受限门与全动作读出的决策损失（227条核验断言）。$\Delta$为全动作减受限门。}
\label{tab:readout}
\begin{tabular}{lrrrrrrrrr}
\toprule
\multirow{2}{*}{方法} & \multicolumn{3}{c}{DeepSeek} & \multicolumn{3}{c}{GLM} & \multicolumn{3}{c}{Qwen3} \\
\cmidrule(lr){2-4}\cmidrule(lr){5-7}\cmidrule(lr){8-10}
 & 受限门 & 全动作 & $\Delta$ & 受限门 & 全动作 & $\Delta$ & 受限门 & 全动作 & $\Delta$ \\
\midrule
\textbf{FinPlanRAG (Ours)} & \textbf{3.656} & \textbf{2.907} & \textbf{$-0.749$} & \textbf{3.524} & \textbf{2.943} & \textbf{$-0.581$} & \textbf{3.621} & \textbf{2.925} & \textbf{$-0.696$} \\
RF & 3.454 & 2.987 & $-0.467$ & 3.454 & 2.987 & $-0.467$ & 3.454 & 2.987 & $-0.467$ \\
LLM-only & 3.674 & 3.445 & $-0.229$ & 3.454 & 3.427 & $-0.026$ & 5.048 & 5.498 & $+0.449$ \\
\bottomrule
\end{tabular}
\end{table}

全动作读出的作用集中在高不确定性区域。融合后验在227条测试断言中有217条选择裁决集、10条选择点答案、0条弃权；平均集合大小为2.12。点答案不是唯一合理的输出，当两个相邻结果的期望损失接近时，返回二元裁决集可以降低误判成本。这一结果说明，前瞻性断言的核心问题不只是预测哪个标签，还包括给定后验应当输出什么。

\section{指导区间内的分层结果}
\label{sec:within}

总体损失打平时，锚点层是否在某个误差层上提供互补信号，是一个独立问题。表\ref{tab:within}报告真实结果落在指引区间内的82条测试断言。FinPlanRAG在三个基座上均低于随机森林：DeepSeek从3.146降到2.780，GLM从3.146降到2.902，Qwen从3.146降到2.732（以第一个随机森林种子为例）。五个随机森林种子中，DeepSeek和Qwen的五个区间全部排除零，GLM有四个区间排除零。

\begin{table}[htbp]
\centering
\zihao{5}
\setlength{\tabcolsep}{4pt}
\caption{指导区间内样本的分层比较（82条测试断言）。差值为FinPlanRAG减随机森林，负值表示FinPlanRAG损失更低。}
\label{tab:within}
\begin{tabular}{lrrrr}
\toprule
基座 & FinPlanRAG & 随机森林 & 损失差 [95\% CI] & 种子一致性 \\
\midrule
DeepSeek & \textbf{2.780} & 3.146 & $-0.366$ $[-0.732,-0.049]$ & 5/5 \\
GLM & \textbf{2.902} & 3.146 & $-0.244$ $[-0.537,-0.024]$ & 4/5 \\
Qwen3 & \textbf{2.732} & 3.146 & $-0.415$ $[-0.732,-0.170]$ & 5/5 \\
\bottomrule
\end{tabular}
\end{table}

这个分层结果说明锚点层在高密度、低极端误差的区域提供了互补信号。分层依据是真实结果，预测时不可知，因此它只能解释误差来源，不能作为部署时的路由规则。总体损失上的差异仍按不显著报告。

\section{A股年报历史回测}
\label{sec:ashare}

表\ref{tab:ashare}给出A股断言的分类与决策结果，DeepSeek和GLM使用470条全样本，Qwen3使用250条分层子样本。上年外推基线把上年增速或上年毛利率直接外推，季度年化基线把三季报年化后外推；完整系统使用锚点、资格判定、时间演化、数值桥和全动作读出，去桥变体关闭数值桥。

\begin{table}[htbp]
\centering
\zihao{-5}
\setlength{\tabcolsep}{3pt}
\caption{A股年报历史回测。Acc.为准确率，Bal.为平衡准确率，Macro为宏F1，Loss为全动作决策损失，日期为披露日期窗口准确率。Qwen3在250条分层子样本上评估，其余方法在470条全样本上评估。}
\label{tab:ashare}
\begin{tabular}{llrrrrr}
\toprule
方法 & 样本 & Acc. & Bal. & Macro & Loss & 日期 \\
\midrule
\multicolumn{7}{l}{\textit{传统基线}} \\
上年外推 & 470 & .491 & .402 & .401 & 5.094 & .780 \\
季度年化 & 470 & .777 & .786 & .785 & 2.270 & .729 \\
\midrule
\multicolumn{7}{l}{\textit{DeepSeek}} \\
\textbf{FinPlanRAG} & 470 & \textbf{.783} & \textbf{.778} & \textbf{.776} & \textbf{2.179} & \textbf{.788} \\
去数值桥 & 470 & .638 & .591 & .590 & 2.421 & .788 \\
\midrule
\multicolumn{7}{l}{\textit{GLM}} \\
\textbf{FinPlanRAG} & 470 & \textbf{.549} & \textbf{.485} & \textbf{.483} & \textbf{3.332} & \textbf{.754} \\
去数值桥 & 470 & .549 & .485 & .483 & 3.332 & .754 \\
\midrule
\multicolumn{7}{l}{\textit{Qwen3}} \\
季度年化 & 250 & .772 & .789 & .789 & 2.324 & .714 \\
\textbf{FinPlanRAG} & 250 & \textbf{.580} & \textbf{.523} & \textbf{.522} & \textbf{4.208} & \textbf{.690} \\
去数值桥 & 250 & .564 & .500 & .499 & 3.864 & .690 \\
\bottomrule
\end{tabular}
\end{table}

DeepSeek基座上的完整系统取得.783的准确率和2.179的决策损失，高于季度年化的.777和2.270，也明显高于去桥变体的.638和2.421。完整系统相对去桥变体的损失差为$-0.243$ $[-0.404,-0.077]$，准确率差为$+0.145$ $[0.109,0.183]$，两个区间均排除零。数值桥把营业收入和归母净利润的增速判断从标签后验推进到可计算的数值区间，是DeepSeek增益的主要来源。GLM基座上完整系统与去桥变体相同，说明该基座的抽取结果没有提供足够的数值锚点，桥接未触发；完整系统也低于季度年化。Qwen3的250条子样本上，完整系统的准确率比去桥变体高$+0.016$ $[0.004,0.032]$，但损失高$+0.344$ $[0.064,0.640]$，说明数值桥改善了部分区间判断，却没有改善整体决策校准。A股结果因此分为两类：DeepSeek上数值桥带来明确增益，GLM和Qwen上抽取质量不足以支撑桥接，标签后验仍是更稳的回退。

表\ref{tab:ashare-numeric}给出数值桥的点估计。DeepSeek完整系统的数值MAE为9.86个百分点，季度年化为7.40个百分点；完整系统的分类损失更低，但数值点估计不如简单年化。这说明两种目标并不等价：数值桥改善的是区间判断和决策损失，不保证点值误差最小。Qwen3完整系统的数值MAE为39.21个百分点，明显高于DeepSeek，说明其抽取器输出的数值字段噪声较大。GLM没有产生可用的数值点估计，数值桥未触发。披露日期窗口准确率见表\ref{tab:ashare}：DeepSeek为.788，GLM为.754，Qwen3为.690；对应的季度年化分别为.729、.729和.714。

\begin{table}[htbp]
\centering
\zihao{5}
\setlength{\tabcolsep}{4pt}
\caption{A股回测的数值桥点估计。MAE为增速或毛利率的百分点误差，MAPE只在真实值非零时计算。GLM没有产生可用的数值点估计，数值桥未触发。}
\label{tab:ashare-numeric}
\begin{tabular}{llrr}
\toprule
方法 & 样本 & 数值MAE & 数值MAPE \\
\midrule
季度年化 & 470 & 7.40 & 459.7 \\
\textbf{FinPlanRAG (DeepSeek)} & 470 & \textbf{9.86} & \textbf{408.7} \\
\textbf{FinPlanRAG (Qwen3)} & 250 & \textbf{39.21} & \textbf{508.7} \\
\bottomrule
\end{tabular}
\end{table}

\section{公开基准的组件级检查}

表\ref{tab:findver-component}给出FinDVer前150条断言的结果。状态条件式获取从两条证据开始，只有在阅读器报告证据不足时才扩展，最多八条；固定预算条件直接给4条或8条；金标条件给出专家标注的证据段落。

\begin{table}[H]
\centering
\zihao{5}
\setlength{\tabcolsep}{4pt}
\caption{FinDVer组件检查（150条断言）。Acc.为断言判定准确率，Macro为宏F1，Recall为证据段落召回率，Para.为平均证据段落数。}
\label{tab:findver-component}
\zihao{-5}
\setlength{\tabcolsep}{2.2pt}
\begin{tabular}{lrrrrrrrrrrrr}
\toprule
\multirow{2}{*}{条件} & \multicolumn{4}{c}{DeepSeek} & \multicolumn{4}{c}{GLM} & \multicolumn{4}{c}{Qwen3} \\
\cmidrule(lr){2-5}\cmidrule(lr){6-9}\cmidrule(lr){10-13}
 & Acc. & Macro & Recall & Para. & Acc. & Macro & Recall & Para. & Acc. & Macro & Recall & Para. \\
\midrule
固定4条 & .767 & .550 & .620 & 4.00 & .760 & .485 & .620 & 4.00 & .733 & .438 & .620 & 4.00 \\
固定8条 & .847 & .570 & .755 & 8.00 & .860 & .573 & .755 & 8.00 & .700 & .426 & .755 & 8.00 \\
金标证据 & .913 & .577 & 1.000 & 2.99 & .933 & .603 & 1.000 & 2.99 & .707 & .441 & 1.000 & 2.99 \\
\textbf{FinPlanRAG 获取} & \textbf{.873} & \textbf{.559} & \textbf{.603} & \textbf{4.63} & \textbf{.913} & \textbf{.592} & \textbf{.628} & \textbf{4.95} & \textbf{.660} & \textbf{.397} & \textbf{.517} & \textbf{4.01} \\
\bottomrule
\end{tabular}
\end{table}

DeepSeek和GLM上，FinPlanRAG的状态条件式获取在约一半段落数下取得高于固定8条的判定准确率：DeepSeek为.873对.847，GLM为.913对.860。证据召回率低于固定8条，说明它没有覆盖全部金标段落，但覆盖了足以判定的段落。Qwen3上的方向相反：固定4条为.733，状态条件式为.660，低于金标条件.707。Qwen3的阅读器在短上下文下更稳定，扩展证据反而引入噪声。这个结果把状态条件式获取的适用条件限定在阅读器能够利用新增证据的基座上。



% ==================== 第7章 消融实验 ====================
\section{消融实验}

本章的消融在锚点分支和最终系统上分别进行。锚点分支的损失只用于比较证据策略，不与第\ref{ch:empirical}章的融合主表直接比较；最终系统的消融使用与主表相同的后验和读出。

\subsection{锚点提示词与时间演化}

文本锚点的作用首先出现在纯LLM分支。旧提示词使用结构化摘要，新提示词使用时点原文、段落编号和逐字引文。DeepSeek上的准确率、平衡准确率和宏F1分别从.445、.307、.242提高到.502、.397、.364。来源落地和逐字引文让模型能够区分“指引区间”“历史实际”和“分析师预测”，减少了把累计增速当作全年增速的错误。

时间演化的作用较小但方向一致。新锚点从静态可靠度切换到按揭示日更新的Beta可靠度，损失从2.908降到2.855；旧锚点从2.899降到2.793。准确率变化不大，演化主要改善概率校准和决策损失，而不是点分类。这个结果符合演化的设计目标：它更新的是锚点类型的可靠性，不是标签本身。

\subsection{资格判定}

表\ref{tab:qual-ablation}比较完整资格判定与接受全部解析锚点。完整判定拒绝空泛、过期、对抗和不合格的锚点，接受率在三个基座上分别为27.6\%、42.7\%和68.3\%；接受全部条件保留所有解析结果。

\begin{table}[htbp]
\centering
\zihao{5}
\setlength{\tabcolsep}{4pt}
\caption{资格判定消融（227条核验断言，锚点分支）。差值为接受全部减完整判定，正值表示完整判定损失更低。}
\label{tab:qual-ablation}
\begin{tabular}{lrrrr}
\toprule
基座 & 完整判定 & 接受全部 & 损失差 [95\% CI] & 接受率 \\
\midrule
DeepSeek & 3.304 & 3.507 & $+0.203$ $[-0.106,0.520]$ & 27.6\% \\
GLM & 3.339 & 3.427 & $+0.088$ $[-0.176,0.361]$ & 42.7\% \\
Qwen3 & 4.599 & 4.652 & $+0.053$ $[-0.141,0.256]$ & 68.3\% \\
\bottomrule
\end{tabular}
\end{table}

三个基座的点估计都支持完整判定，DeepSeek的增益最大，Qwen3的增益最小。三个配对区间都包含零，因此资格判定的作用按方向性证据报告。Qwen3的接受率明显高于另外两个基座，说明其抽取器产生了更多边缘锚点，资格判定拒绝的比例更高，但剩余锚点的区分力有限。

\subsection{融合权重}

表\ref{tab:weight}给出验证集上不同融合权重的决策损失。权重$w$是锚点后验在融合后验中的占比，$w=0$对应纯随机森林，$w=1$对应纯锚点后验。DeepSeek在$w=0.2$取得最低验证损失，GLM和Qwen3在$w=0.1$取得最低验证损失；测试集使用验证集选择的权重，不参与调参。

\begin{table}[htbp]
\centering
\zihao{5}
\setlength{\tabcolsep}{4pt}
\caption{融合权重的验证集损失（185条验证断言）。粗体为验证集选择。}
\label{tab:weight}
\begin{tabular}{lrrrrrrr}
\toprule
基座 & 0.0 & 0.05 & 0.1 & 0.2 & 0.3 & 0.5 & 1.0 \\
\midrule
DeepSeek & 3.341 & 3.308 & 3.265 & \textbf{3.232} & 3.341 & 3.503 & 3.859 \\
GLM & 3.341 & 3.211 & \textbf{3.200} & 3.286 & 3.362 & 3.416 & 3.708 \\
Qwen3 & 3.341 & 3.362 & \textbf{3.330} & 3.686 & 3.859 & 4.162 & 6.519 \\
\bottomrule
\end{tabular}
\end{table}

三个基座的最优权重都偏向随机森林，DeepSeek为0.2，GLM和Qwen3为0.1。锚点后验单独使用时损失明显更高，Qwen3在$w=1$时达到6.519。融合的价值来自补充信号，而不是替代结构化后验。

\subsection{类型库的验证集扩展}

类型库扩展是算法的一个训练步骤，不是独立的预测模块。算法\ref{alg:typeexpand}先在验证集上挖掘高损失断言，由LLM提出候选类型，再在选择集上重抽锚点并计算损失；只有验证损失下降超过$\delta$的候选才进入类型库。接受后的类型库冻结，测试集只评估一次。


DeepSeek上的候选类型库包含“指引中点与当前季度收入之比”和“财年周数效应”两个类型。验证选择集上，基础类型库的损失为3.943，候选类型库为3.792，下降0.151，超过阈值，候选被接受。提示词单独修改和类型加提示词组合两个候选没有通过验证门控。冻结类型库后，测试集上的准确率从.555提高到.577，宏F1从.513提高到.530，损失从2.934降到2.907；配对损失差为$-0.026$ $[-0.159,0.097]$，区间包含零。

表\ref{tab:typeevolve}汇总这一过程。类型库扩展的测试增益是小的，但方向与验证集一致，且没有使用测试标签。GLM和Qwen3没有候选通过验证门控，保留基础类型库，因此主表中只有DeepSeek使用扩展类型库。

\begin{table}[htbp]
\centering
\zihao{5}
\setlength{\tabcolsep}{4pt}
\caption{DeepSeek类型库扩展（227条核验断言）。Loss为全动作决策损失。}
\label{tab:typeevolve}
\begin{tabular}{lrrrr}
\toprule
系统 & Acc. & Bal. & Macro & Loss \\
\midrule
基础类型库 & .555 & .511 & .513 & 2.934 \\
扩展类型库 & .577 & .526 & .530 & 2.907 \\
差值 & $+.022$ & $+.015$ & $+.017$ & $-0.026$ $[-0.159,0.097]$ \\
\bottomrule
\end{tabular}
\end{table}

\subsection{去掉锚点结构的对照}
\label{sec:no-anchor}

前面的消融检验的是组件内部，例如资格判定是否比接受全部更好。一个更根本的问题是：如果去掉三层锚点结构，只保留时点证据、LLM标签分数、温度缩放和全动作读出，损失会差多少。这个对照直接检验锚点结构本身的必要性。

本文从同一批锚点抽取响应中只保留标签分数，去掉全部锚点对象、资格判定和Beta可靠性，在验证集上拟合温度参数，再在测试集上使用同一个全动作读出。这个对照可以看作FinAbstain式选择性预测在本任务上的近似：时点证据、概率输出、温度缩放和损失感知弃权都在，但没有多类别资格判定和类型可靠性演化。表\ref{tab:noanchor}给出结果。去掉锚点结构后，DeepSeek、GLM和Qwen的损失分别上升0.379、0.317和0.899，准确率分别下降5.3、8.8和13.7个百分点。时点证据和选择性输出仍然存在，性能却明显低于完整系统，说明增益不能只归因于“时点检索加弃权”这一一般性组合。

\begin{table}[htbp]
\centering
\zihao{5}
\setlength{\tabcolsep}{4pt}
\caption{去掉三层锚点结构的架构级对照（227条核验断言）。温度在验证集上选择，测试集只评估一次。}
\label{tab:noanchor}
\begin{tabular}{lrrrrrrrrr}
\toprule
\multirow{2}{*}{方法} & \multicolumn{3}{c}{DeepSeek} & \multicolumn{3}{c}{GLM} & \multicolumn{3}{c}{Qwen3} \\
\cmidrule(lr){2-4}\cmidrule(lr){5-7}\cmidrule(lr){8-10}
 & Acc. & Macro & Loss & Acc. & Macro & Loss & Acc. & Macro & Loss \\
\midrule
无锚点结构 & .524 & .398 & 3.286 & .463 & .347 & 3.260 & .405 & .257 & 3.824 \\
\textbf{FinPlanRAG} & \textbf{.577} & \textbf{.530} & \textbf{2.907} & \textbf{.551} & \textbf{.504} & \textbf{2.943} & \textbf{.542} & \textbf{.501} & \textbf{2.925} \\
\bottomrule
\end{tabular}
\end{table}

这个对照也解释了为什么本文把贡献放在证据资格判定和类型可靠性演化上，而不是放在“选择性预测”上。选择性预测本身是一个一般性工具，FinAbstain等工作已经把它用于金融预测；本文的增量在于锚点结构把证据拆成可判定、可追溯、可按类型更新的对象。去掉这层结构，系统仍然可以选择弃权或裁决集，但失去了区分证据失效模式的能力。

\subsection{数值桥}

数值桥的作用在A股回测上最清楚。DeepSeek完整系统的损失为2.179，去桥变体为2.421，差值为$-0.242$；准确率从.638提高到.783，宏F1从.590提高到.776。桥接把“增速可能高于区间”的标签判断转化为“按当前季度年化，全年增速约为多少”的数值计算，减少了纯标签后验的过度平滑。

桥接的收益依赖抽取质量。GLM完整系统与去桥变体完全相同，说明该基座没有产生足够的数值锚点，桥接未触发；Qwen3的结果见第\ref{ch:empirical}章。数值桥不是无条件增益，它要求抽取器输出可核验的数值字段、单位和期间。对抽取质量较弱的基座，标签后验仍是更稳的回退。

\subsection{消融结论}

消融结果可以归纳为四点。第一，去掉三层锚点结构后，三个基座的损失分别上升0.379、0.317和0.899，说明锚点结构本身提供了不可替代的增量，而不是只有选择性输出在起作用。第二，输出形式和资格判定是稳定的组件：全动作读出在三个基座上显著降低损失，资格判定方向一致但区间包含零。第三，类型库扩展和数值桥是条件增益：它们在DeepSeek上有效，在抽取质量较弱的基座上不触发或收益有限。第四，融合权重偏向随机森林，锚点层提供的是补充信号，而不是替代结构化后验。这个结论与主表的定位一致：框架的优势来自证据状态、决策接口和可审计性，而不是单一组件的普遍提升。


\section{外部迁移与误差分析}

表\ref{tab:findver}报告FinDVer testmini默认阅读器结果。DeepSeek Flash上规划器加回退取得0.701准确率，平均8.76段；固定16段0.660，固定8段0.623。GLM 5.3 Flash上该分支取得0.753，平均9.19段；固定16段0.743，固定8段0.680。Qwen3-8B上该分支取得0.649，平均6.13段；固定16段0.636，固定8段0.629。分支在三个基座上均高于固定16段阅读器，平均少用6.9至9.9段。官方思维链阅读器在DeepSeek Flash和GLM 5.3 Flash上分别取得0.817和0.819（金证据），十段检索版本均为0.787；这些使用不同阅读器和证据接口，是上界参考，不是同条件基线。FinDVer是迁移集，分支未在其上调参。

义务检索分支的主要错误是证据缺失和算术错误。DocFinQA的段落块可能把表格行拆开，导致阅读器找不到分子或分母；规划器的\texttt{missing}字段可以部分修复，但当缺失项不在BM25候选的前32条中时无法恢复。Qwen3-8B在DocFinQA上的精确匹配接近地板，主要原因是回答包含解释性文字，而不是数值本身错误。这些错误指向两个改进方向：更细的表格块边界和更严格的短答案格式约束。

AnchorEvolve-LLM的主要误差来自类别偏置。DeepSeek Flash较多地预测高于区间，纯LLM后验的平衡准确率明显偏低；GLM 5.3 Flash与Qwen3-8B的宏F1更高，但概率质量仍不足以单独承担低损失决策。融合后验在三个基座的准确率、平衡准确率、宏F1和决策损失上均为最高，说明LLM锚点信号与结构化特征并不重复。硬锚点触发较少，当前收益主要来自软锚点的重新加权；全动作读出把多数断言送到二元裁决集，平均集合大小约为2.1，说明后验仍偏保守。共享校准随机森林是强诊断基线，其总体损失与FinPlanRAG接近，但在指导区间内样本上显著更差。

路由器的错误集中在语义边界。DeepSeek Flash和GLM 5.3 Flash把“未来支付金额”“预期股利”等确定问题判为前瞻性断言；这些问题的答案在文件中已经给出，但文本包含未来时态。Qwen3-8B的错误更多，说明小模型的语义边界判断不稳定。对于部署，路由错误比分支错误更容易诊断，因为路由输出带有置信度和理由；系统可以在低置信度时请求人工复核，或回退到保守契约。本文不进一步设计这一回退机制，但结果为后续工作提供了接口。

\begin{table}[H]
\centering
\zihao{5}
\setlength{\tabcolsep}{4pt}
\caption{FinDVer testmini迁移结果（默认阅读器）。固定4段、固定8段和规划器单独运行的结果见正文。}
\label{tab:findver}
\begin{tabular}{llrrr}
\toprule
基座 & 条件 & 准确率 & 宏F1 & 平均段落 \\
\midrule
DeepSeek Flash & 固定16段 & .660 & .735 & 16.00 \\
 & 规划器+回退 & .701 & .759 & 8.76 \\
GLM 5.3 Flash & 固定16段 & .743 & .792 & 16.00 \\
 & 规划器+回退 & .753 & .789 & 9.19 \\
Qwen3-8B & 固定16段 & .636 & .673 & 16.00 \\
 & 规划器+回退 & .649 & .663 & 6.13 \\
\bottomrule
\end{tabular}
\end{table}

% ==================== 第7章 ====================
\chapter{讨论与结论}
\label{ch:conclusion}

\section{讨论}

路由实验回答第一个研究问题。DeepSeek Flash和GLM 5.3 Flash分别达到0.995和0.990的路由准确率，混合效用为0.403和0.419，后者恢复oracle的99.3\%。错路由主要发生在包含未来时态、但答案已经在文件中给出的确定问题。Qwen3-8B的路由准确率降至0.867，混合效用0.263，低于Always-Anchor的0.298。这个对比说明契约识别不是简单的时态匹配：较强基座能够利用查询与截止时点的关系，较小基座则容易把语义边界压缩为表面模式。路由实验因此支持H1在较强基座上成立，同时给出基座能力的边界。

DocFinQA和FinDVer回答第二个研究问题。DeepSeek Flash上，义务检索分支取得0.457精确匹配、平均10.69段，固定16段阅读器取得0.390、平均16段；GLM 5.3 Flash上分别为0.497对0.480，平均10.31段对16段。FinDVer上三个基座分别为0.701、0.753和0.649，均高于固定16段的0.660、0.743和0.636。四个具名RAG适配在DocFinQA上均低于义务检索分支，分支的答案覆盖率0.770和0.810也高于固定预算。规划器单独运行的成本更低，但在GLM 5.3 Flash上留下小缺口，置信度回退在24\%的问题上触发。实验支持H2：类型化义务与覆盖驱动取证在长文档上以更低证据成本取得更高准确率，并可迁移到FinDVer；收益来自证据策略，而不是阅读器或评分口径。

收入指引实验回答第三个研究问题，也给出最需要谨慎解释的结果。纯LLM锚点后验在三个基座上的准确率为0.445、0.454和0.427，决策损失为4.035、3.903和3.586，说明锚点包含非随机信号，但类别偏置和概率校准不足以直接承担低损失决策。融合后验把准确率提高到0.524、0.537和0.537，决策损失降到3.604、3.648和3.577，低于六个具名基线；在GLM 5.3 Flash和Qwen3-8B上，准确率也高于最强具名基线。纯RF诊断取得0.533、0.491、0.486和3.454，是当前最强的数值决策头。这个结果不支持“LLM锚点后验优于强RF”，但支持“语义状态与校准读出分工”：LLM负责路由、阅读和锚点状态，RF负责数值风险控制。H3应表述为融合读出相对具名基线降低决策损失，而不是LLM单独胜过RF。

消融实验进一步说明锚点层的作用与边界。去掉全部锚点后，LLM后验的宏F1从0.233、0.286、0.339降到0.213、0.211、0.272；取消硬锚点掩码后指标几乎不变，说明硬锚点在收入指引上很少触发；逐类去除锚点没有一致赢家，说明锚点类型的作用依赖基座和数据分布。输出形式对照显示，强制点答案的平均损失为4.693，固定三桶全集为4.000，损失感知门为3.454；风险收益来自集合与弃权，而不是更强的点分类器。误差分析显示，DocFinQA的错误主要来自表格块边界与算术，Qwen3-8B的错误主要来自解释性输出，AnchorEvolve的错误主要来自类别偏置，路由错误集中在语义边界。这些结果把改进方向限定在证据表示、锚点过滤和路由信号上，而不是盲目增加模型规模。

实验比较的是同基座适配，不是原系统的未修改端到端运行；标签来自XBRL核验而非人工金标准；混合负载由DocFinQA与收入指引样本构成。这些选择保证了阅读器、证据接口和评分口径可比，也意味着结论适用于给定的数据、基座和评价口径。它们不改变三个研究问题的检验结果，但界定了结果的外推范围。


\section{局限与未来研究}

本文的证据边界集中在三点。第一，基准覆盖仍有限。收入指引主基准只覆盖美国公司的季度营业收入指引；A股回测覆盖40家公司和三个财年，但只到年报、营收、利润、毛利率和披露日期；产能、利用率、良率和排产等更细的数值桥只有模板和前瞻性案例，没有历史端到端基准。混合负载由DocFinQA和收入指引样本构造，真实用户查询的契约边界可能更复杂。

第二，标签来自XBRL核验或披露文件回退，不是人工金标准；A股数值来自公开财务报表，口径差异和重述仍会带来噪声。Qwen3-8B在路由和数值抽取上明显弱于另外两个基座，GLM的数值桥没有触发，说明框架效果依赖基座的契约识别和字段抽取能力。

第三，类型库扩展只在DeepSeek上完成一轮，GLM和Qwen没有候选通过验证门控；全动作读出在总体损失上稳健，但融合与随机森林的总体差异仍不显著。后续工作将补产能与排产的历史基准、更难的路由边界样本和人工审计，并把类型库扩展扩展到多轮和跨基座迁移。

\section{研究结论}

本文提出FinPlanRAG，把金融混合查询的回答形式选择表述为查询契约路由，并以义务检索和AnchorEvolve分别处理确定信息与前瞻性断言。路由器在检索与生成之前判断事实是否已在截止时点存在；义务检索围绕类型化义务形成闭合答案；AnchorEvolve以L0稳定关系、L1行业模式和L2实例锚点维护可更新、可追溯的前瞻性证据状态。校准读出在此基础上生成点答案、可信集或弃权，用于控制最终决策风险。

第一项贡献是查询契约路由。现有问答、断言验证和选择性预测分别固定输出接口，本文把“事实是否已在截止时点存在”置于检索与生成之前，用一次路由决定闭合答案或可校准判断。路由实验支持这一贡献：DeepSeek Flash和GLM 5.3 Flash分别达到0.995和0.990的路由准确率，混合效用接近oracle，并显著高于Always-RAG和Always-Anchor；Qwen3-8B的结果则说明路由能力依赖基座。契约路由的价值不只是提高一个分类指标，而是把混合负载的分析单位从数据集任务改为查询契约，使路由错误和分支错误可以分开诊断。

第二项贡献是义务检索与AnchorEvolve的两条证据分支。义务检索把确定问题绑定为类型化义务，在时点语料上按覆盖取证，由阅读器返回带证据编号的闭合答案；AnchorEvolve把前瞻性断言映射到有限结果空间，用三层锚点、五类资格判定和时点信念更新保留证据的作用域、可靠性和时效。DocFinQA和FinDVer实验支持义务检索的效率主张；收入指引实验和锚点消融支持AnchorEvolve的语义主张。特别是去除全部锚点后LLM后验宏F1下降，说明锚点不是装饰，而是改变了模型判断的中间状态。

第三项贡献是语义状态与校准读出的分工，以及统一的多基准评估。本文在三个数据集、三套基座和六个具名基线上同时检验路由、阅读、锚点推断、后验融合、校准选择和外部迁移，并把分类、校准、证据成本、决策损失和路由效用放在同一口径下报告。需要明确，纯RF是当前最强的数值决策头，融合后验在部分基座上取得更高普通准确率，但决策损失仍高于纯RF；因此本文不主张LLM锚点后验单独胜过强RF。本文主张的是：LLM在路由、阅读和语义状态构造上不可替代，校准读出负责把这些状态转换为低风险决策，二者应当分开评估。

本文的原创性贡献可以概括为三点。第一，把输出契约从固定接口变成查询层的路由变量，使确定问答与前瞻性裁决在同一时点约束下被统一处理。第二，把前瞻性证据从单一标签变成可审计的状态，使证据来源、作用域、资格理由和信念变化可以被逐条检查。第三，把语义判断与风险校准分开，使LLM的语义能力、统计模型的风险控制能力和路由器的契约识别能力可以分别检验、替换和组合。这三点不依赖特定基座或特定行业，构成了本文相对于现有金融RAG和断言验证工作的主要区别。

从金融实践看，框架把答案、证据和风险放在同一条审计链上。分析师对确定问题获得带证据编号的闭合答案，对前瞻性断言获得点判断、可信集或弃权，并可以查看锚点来源和资格理由。机构可以按自身风险偏好替换损失参数和校准后端，而不改变路由与证据状态的结构。生产里程碑、产能投放、监管审批和合同履约等任务与收入指引具有相似的时点结构，本文结果为这些场景提供可复用的接口。

本文的结论有明确边界。收入指引实验采用三桶结果空间，通用的时机与阶段状态仍需在更多前瞻性事件上评估；混合负载由DocFinQA与收入指引样本构成，真实用户查询的语义边界可能更复杂；标签来自XBRL核验而非人工金标准。后续工作将围绕更多前瞻性事件、更稳定的边界路由、更窄的可信集和更严格的人工复核展开。由于路由、锚点抽取和输出读出彼此分离，这些改进可以独立评估，而不需要改变FinPlanRAG的整体接口。


\clearpage
\phantomsection
\addcontentsline{toc}{chapter}{参考文献}
{\zihao{5}\setlength{\baselineskip}{16pt}\begin{thebibliography}{33}
\setlength{\bibsep}{0pt}
\setlength{\itemsep}{0pt}
\setlength{\parsep}{0pt}
\setlength{\bibhang}{2em}

\bibitem[Angelopoulos and Bates(2021)]{angelopoulos2021gentle}
Angelopoulos, A. N., and Bates, S., 2021, ``A Gentle Introduction to Conformal Prediction and Distribution-Free Uncertainty Quantification'', arXiv:2107.07511.

\bibitem[Asai et al.(2024)]{asai2024selfrag}
Asai, A., Wu, Z., Wang, Y., Sil, A., and Hajishirzi, H., 2024, ``Self-RAG: Learning to Retrieve, Generate, and Critique through Self-Reflection'', in \emph{International Conference on Learning Representations}.

\bibitem[Brodersen et al.(2010)]{brodersen2010balanced}
Brodersen, K. H., Ong, C. S., Stephan, K. E., and Buhmann, J. M., 2010, ``The Balanced Accuracy and Its Posterior Distribution'', in \emph{Proceedings of the 20th International Conference on Pattern Recognition}, pp. 3121--3124.

\bibitem[Chen et al.(2026)]{chen2026hcrag}
Chen, S., Tan, H., Li, Y., and Liang, J., 2026, ``HC-RAG: Evidence-Centric Retrieval-Augmented Generation over Heterogeneous Financial Filings'', arXiv:2608.12335.

\bibitem[Chen et al.(2021)]{chen2021finqa}
Chen, Z., Chen, W., Smiley, C., Shah, S., Borova, I., Langdon, D., Moussa, R., Beane, M., Huang, T.-H., Routledge, B., and Wang, W. Y., 2021, ``FinQA: A Dataset of Numerical Reasoning over Financial Data'', in \emph{Proceedings of the 2021 Conference on Empirical Methods in Natural Language Processing}, pp. 3697--3711.

\bibitem[Chi et al.(2026)]{chi2026finsagent}
Chi, J., Tai, Z., Wu, H., Kwok, T. S. T., He, H., Liao, Z., Xiao, B., Jiang, C., Lei, J., Huang, J., Lu, P., Li, M., Ma, L., Wu, Y., Lyu, S., Tian, J., Li, Y., Ma, Y., Guan, S., Hu, D., Cui, Y., Zhou, L., Ding, L., and Wang, X., 2026, ``FinSAgent: Corpus-Aligned Multi-Agent RAG Framework for Evidence-Grounded SEC Filing Question Answering'', arXiv:2607.18102.

\bibitem[Chow(1970)]{chow1970optimum}
Chow, C. K., 1970, ``On Optimum Recognition Error and Reject Tradeoff'', \emph{IEEE Transactions on Information Theory}, Vol. 16, No. 1, pp. 41--46.

\bibitem[Elkan(2001)]{elkan2001cost}
Elkan, C., 2001, ``The Foundations of Cost-Sensitive Learning'', in \emph{Proceedings of the 17th International Joint Conference on Artificial Intelligence}, pp. 973--978.

\bibitem[Geifman and El-Yaniv(2017)]{geifman2017selective}
Geifman, Y., and El-Yaniv, R., 2017, ``Selective Classification for Deep Neural Networks'', in \emph{Advances in Neural Information Processing Systems}, Vol. 30.

\bibitem[Gibbs and Cand\`es(2021)]{gibbs2021adaptive}
Gibbs, I., and Cand\`es, E. J., 2021, ``Adaptive Conformal Inference Under Distribution Shift'', in \emph{Advances in Neural Information Processing Systems}, Vol. 34, pp. 1660--1672.

\bibitem[Guo et al.(2017)]{guo2017calibration}
Guo, C., Pleiss, G., Sun, Y., and Weinberger, K. Q., 2017, ``On Calibration of Modern Neural Networks'', in \emph{Proceedings of the 34th International Conference on Machine Learning}, Vol. 70, pp. 1321--1330.

\bibitem[Hirst et al.(2008)]{hirst2008management}
Hirst, D. E., Koonce, L., and Venkataraman, S., 2008, ``Management Earnings Forecasts: A Review and Framework'', \emph{Accounting Horizons}, Vol. 22, No. 3, pp. 315--338.

\bibitem[Islam et al.(2023)]{islam2023financebench}
Islam, P., Kannappan, A., Kiela, D., Qian, R., Scherrer, N., and Vidgen, B., 2023, ``FinanceBench: A New Benchmark for Financial Question Answering'', arXiv:2311.11944.

\bibitem[Jeong et al.(2024)]{jeong2024adaptiverag}
Jeong, S., Baek, J., Cho, S., Hwang, S. J., and Park, J. C., 2024, ``Adaptive-RAG: Learning to Adapt Retrieval-Augmented Large Language Models through Question Complexity'', in \emph{Proceedings of the 2024 Conference of the North American Chapter of the Association for Computational Linguistics}, pp. 7036--7050.

\bibitem[Jiang et al.(2026)]{jiang2026finrate}
Jiang, Y., Chen, J., Makri, E., Chen, J., Li, P., Maatouk, A., Tassiulas, L., Brenner, E., Xiang, B., and Ying, R., 2026, ``Fin-RATE: A Real-world Financial Analytics and Tracking Evaluation Benchmark for LLMs on SEC Filings'', in \emph{Proceedings of the 32nd ACM SIGKDD Conference on Knowledge Discovery and Data Mining}, pp. 9183--9194.

\bibitem[Jiang et al.(2023)]{jiang2023flare}
Jiang, Z., Xu, F., Gao, L., Sun, Z., Liu, Q., Dwivedi-Yu, J., Yang, Y., Callan, J., and Neubig, G., 2023, ``Active Retrieval Augmented Generation'', in \emph{Proceedings of the 2023 Conference on Empirical Methods in Natural Language Processing}, pp. 7969--7992.

\bibitem[Lewis et al.(2020)]{lewis2020rag}
Lewis, P., Perez, E., Piktus, A., Petroni, F., Karpukhin, V., Goyal, N., Kuttler, H., Lewis, M., Yih, W.-T., Rockt\"aschel, T., Riedel, S., and Kiela, D., 2020, ``Retrieval-Augmented Generation for Knowledge-Intensive NLP Tasks'', in \emph{Advances in Neural Information Processing Systems}, Vol. 33, pp. 9459--9474.

\bibitem[L\'opez de Prado(2018)]{lopezdeprado2018advances}
L\'opez de Prado, M., 2018, \emph{Advances in Financial Machine Learning}, Wiley.

\bibitem[Manski(2003)]{manski2003partial}
Manski, C. F., 2003, \emph{Partial Identification of Probability Distributions}, Springer.

\bibitem[Mansouri et al.(2026)]{mansouri2026finrank}
Mansouri, S., Saad, D., Wahrenburg, M., Weissel, M., and Woebbeking, F., 2026, ``FinRank: An Evidence-Grounded Benchmark for Financial Question Answering and Retrieval over SEC Filings'', arXiv:2608.07400.

\bibitem[Qwen Team(2025)]{qwen2025qwen3}
Qwen Team, 2025, ``Qwen3 Technical Report'', arXiv:2505.09388.

\bibitem[Rangapur et al.(2025)]{rangapur2025finfact}
Rangapur, A., Wang, H., Jian, L., and Shu, K., 2025, ``Fin-Fact: A Benchmark Dataset for Multimodal Financial Fact-Checking and Explanation Generation'', in \emph{Companion Proceedings of the ACM Web Conference 2025}, pp. 785--788.

\bibitem[Reddy et al.(2024)]{reddy2024docfinqa}
Reddy, V., Koncel-Kedziorski, R., Lai, V. D., Krumdick, M., Lovering, C., and Tanner, C., 2024, ``DocFinQA: A Long-Context Financial Reasoning Dataset'', in \emph{Proceedings of the 62nd Annual Meeting of the Association for Computational Linguistics (Volume 2: Short Papers)}, pp. 445--458.

\bibitem[Romano et al.(2020)]{romano2020classification}
Romano, Y., Sesia, M., and Cand\`es, E. J., 2020, ``Classification with Valid and Adaptive Coverage'', in \emph{Advances in Neural Information Processing Systems}, Vol. 33, pp. 3581--3591.

\bibitem[Shinn et al.(2023)]{shinn2023reflexion}
Shinn, N., Cassano, F., Berman, A., Gopinath, A., Narasimhan, K., and Yao, S., 2023, ``Reflexion: Language Agents with Verbal Reinforcement Learning'', in \emph{Advances in Neural Information Processing Systems}, Vol. 36.

\bibitem[Sokolova and Lapalme(2009)]{sokolova2009performance}
Sokolova, M., and Lapalme, G., 2009, ``A Systematic Analysis of Performance Measures for Classification Tasks'', \emph{Information Processing and Management}, Vol. 45, No. 4, pp. 427--437.

\bibitem[Vovk et al.(2005)]{vovk2005algorithmic}
Vovk, V., Gammerman, A., and Shafer, G., 2005, \emph{Algorithmic Learning in a Random World}, Springer.

\bibitem[Xie et al.(2024)]{xie2024finben}
Xie, Q., Han, W., Chen, Z., Xiang, R., Zhang, X., He, Y., Xiao, M., Li, D., Dai, Y., Feng, D., Xu, Y., Kang, H., Kuang, Z., Yuan, C., Yang, K., Luo, Z., Zhang, T., Liu, Z., Xiong, G., Deng, Z., Jiang, Y., Yao, Z., Li, H., Yu, Y., Hu, G., Huang, J., Liu, X.-Y., Lopez-Lira, A., Wang, B., Lai, Y., Wang, H., Peng, M., Ananiadou, S., and Huang, J., 2024, ``FinBen: A Holistic Financial Benchmark for Large Language Models'', in \emph{Advances in Neural Information Processing Systems}, Vol. 37.

\bibitem[Yan et al.(2024)]{yan2024crag}
Yan, S.-Q., Gu, J.-C., Zhu, Y., and Ling, Z.-H., 2024, ``Corrective Retrieval Augmented Generation'', arXiv:2401.15884.

\bibitem[Yao et al.(2023)]{yao2023react}
Yao, S., Zhao, J., Yu, D., Du, N., Shafran, I., Narasimhan, K., and Cao, Y., 2023, ``ReAct: Synergizing Reasoning and Acting in Language Models'', in \emph{International Conference on Learning Representations}.

\bibitem[Zhao et al.(2024)]{zhao2024findver}
Zhao, Y., Long, Y., Jiang, T., Wang, C., Chen, W., Liu, H., Tang, X., Zhang, Y., Zhao, C., and Cohan, A., 2024, ``FinDVer: Explainable Claim Verification over Long and Hybrid-Content Financial Documents'', in \emph{Proceedings of the 2024 Conference on Empirical Methods in Natural Language Processing}, pp. 14739--14752.

\bibitem[Zhou et al.(2026)]{zhou2026fincards}
Zhou, Y., Zhang, F., Chen, Y., Zhang, H., Nakov, P., and Xie, Z., 2026, ``FinCARDS: Card-Based Analyst Reranking for Financial Document Question Answering'', in \emph{Findings of the Association for Computational Linguistics: ACL 2026}, pp. 24836--24852.

\bibitem[Zhu et al.(2021)]{zhu2021tatqa}
Zhu, F., Lei, W., Huang, Y., Wang, C., Zhang, S., Lv, J., Feng, F., and Chua, T.-S., 2021, ``TAT-QA: A Question Answering Benchmark on a Hybrid of Tabular and Textual Content in Finance'', in \emph{Proceedings of the 59th Annual Meeting of the Association for Computational Linguistics}, pp. 3277--3287.

\end{thebibliography}
}

% ==================== 附录 ====================
\appendix
\chapter{完整指标表}
\label{app:fullmetrics}

表\ref{tab:vague-ds}至表\ref{tab:vague-qwen}报告三个基座上的完整十一项指标。主文表\ref{tab:vague}和表\ref{tab:vague-qwen-main}并列报告AnchorEvolve-LLM、包含LLM信号的AnchorEvolve-Fusion以及共享校准RF诊断，以区分语义状态、融合后验与非LLM风险控制读出。完整指标中的AUROC和AUPRC衡量排序质量，ECE和Brier衡量概率校准，Sel@50和AURC衡量选择性预测，决策损失衡量非对称风险。

\begin{table}[htbp]
\centering
\zihao{5}
\setlength{\tabcolsep}{2pt}
\caption{DeepSeek基座上的模糊断言结果（227条XBRL核验断言）。Loss越低越好；Sel@50为50\%覆盖率下的选择性准确率，AURC越低越好。}
\label{tab:vague-ds}
\resizebox{\linewidth}{!}{%
\begin{tabular}{lrrrrrrrrrrr}
\toprule
方法 & Acc & Bal & Macro & MCC & AUROC & AUPRC & ECE & Brier & Loss & Sel@50 & AURC \\
\midrule
FinDVer & \textbf{.551} & .461 & .464 & \textbf{.233} & .674 & .494 & .092 & .580 & 4.678 & .675 & .337 \\
Fin-Fact & .537 & .446 & .450 & .201 & .650 & .475 & .075 & .590 & 4.705 & .675 & .338 \\
Self-RAG & .546 & .451 & .456 & .216 & .670 & .491 & .102 & .580 & 4.564 & \textbf{.684} & \textbf{.333} \\
CRAG & .529 & .423 & .418 & .177 & .643 & .473 & .097 & .593 & 4.802 & .658 & .349 \\
FLARE & .511 & .417 & .411 & .164 & .634 & .467 & .130 & .610 & 4.881 & .596 & .362 \\
Adaptive-RAG & .493 & .406 & .400 & .143 & .584 & .425 & .147 & .634 & 5.093 & .579 & .415 \\
AnchorEvolve-LLM & .445 & .307 & .242 & -0.075 & .465 & .312 & .129 & .666 & 4.035 & .509 & .529 \\
FinPlanRAG（AnchorEvolve-Fusion, $w=0.1$） & .524 & .471 & .468 & .201 & .703 & .503 & \textbf{.033} & \textbf{.564} & 3.604 & .632 & .383 \\
共享校准RF（非LLM诊断） & .533 & \textbf{.491} & \textbf{.486} & .229 & \textbf{.706} & \textbf{.505} & .033 & .566 & \textbf{3.454} & .649 & .383 \\
\bottomrule
\end{tabular}
}
\end{table}


\begin{table}[htbp]
\centering
\zihao{5}
\setlength{\tabcolsep}{2pt}
\caption{GLM 5.3基座上的模糊断言结果（227条XBRL核验断言）。Loss越低越好；Sel@50为50\%覆盖率下的选择性准确率，AURC越低越好。}
\label{tab:vague-glm}
\resizebox{\linewidth}{!}{%
\begin{tabular}{lrrrrrrrrrrr}
\toprule
方法 & Acc & Bal & Macro & MCC & AUROC & AUPRC & ECE & Brier & Loss & Sel@50 & AURC \\
\midrule
FinDVer & .498 & .439 & .442 & .178 & .642 & .439 & .103 & .627 & 5.233 & .605 & .418 \\
Fin-Fact & .454 & .398 & .402 & .089 & .603 & .430 & .178 & .670 & 5.339 & .579 & .427 \\
Self-RAG & .454 & .406 & .404 & .113 & .617 & .434 & .164 & .637 & 5.163 & .588 & .411 \\
CRAG & .515 & .431 & .435 & .173 & .607 & .418 & .148 & .657 & 4.846 & .570 & .441 \\
FLARE & .471 & .388 & .387 & .110 & .583 & .389 & .192 & .705 & 5.286 & .526 & .518 \\
Adaptive-RAG & .427 & .372 & .376 & .059 & .580 & .382 & .227 & .738 & 5.630 & .482 & .542 \\
AnchorEvolve-LLM & .454 & .329 & .295 & -0.004 & .498 & .339 & .145 & .671 & 3.903 & .430 & .552 \\
FinPlanRAG（AnchorEvolve-Fusion, $w=0.1$） & \textbf{.537} & .489 & \textbf{.487} & .225 & .704 & .504 & \textbf{.032} & \textbf{.562} & 3.648 & .640 & \textbf{.376} \\
共享校准RF（非LLM诊断） & .533 & \textbf{.491} & .486 & \textbf{.229} & \textbf{.706} & \textbf{.505} & .033 & .566 & \textbf{3.454} & \textbf{.649} & .383 \\
\bottomrule
\end{tabular}
}
\end{table}


\begin{table}[htbp]
\centering
\zihao{5}
\setlength{\tabcolsep}{2pt}
\caption{Qwen3-8B基座上的模糊断言结果（227条XBRL核验断言）。Loss越低越好；Sel@50为50\%覆盖率下的选择性准确率，AURC越低越好。}
\label{tab:vague-qwen}
\resizebox{\linewidth}{!}{%
\begin{tabular}{lrrrrrrrrrrr}
\toprule
方法 & Acc & Bal & Macro & MCC & AUROC & AUPRC & ECE & Brier & Loss & Sel@50 & AURC \\
\midrule
FinDVer-style & .445 & .422 & .406 & .133 & .564 & .375 & .555 & 1.110 & 5.524 & .456 & .535 \\
Fin-Fact-style & .511 & .477 & .464 & \textbf{.237} & .609 & .404 & .489 & .978 & 4.846 & .491 & .467 \\
Self-RAG-style & .489 & .447 & .437 & .196 & .583 & .394 & .489 & .990 & 5.093 & .482 & .506 \\
CRAG-style & .454 & .418 & .399 & .152 & .566 & .371 & .547 & 1.093 & 5.366 & .456 & .534 \\
FLARE-style & .326 & .298 & .212 & -0.107 & .475 & .322 & .672 & 1.340 & 6.740 & .316 & .664 \\
Adaptive-RAG-style & .344 & .324 & .290 & -0.027 & .492 & .330 & .654 & 1.308 & 6.485 & .298 & .663 \\
AnchorEvolve-LLM & .427 & .353 & .355 & .027 & .562 & .399 & .163 & .647 & 3.586 & .561 & .435 \\
FinPlanRAG（AnchorEvolve-Fusion, $w=0.1$） & \textbf{.537} & .486 & .485 & .227 & \textbf{.707} & \textbf{.509} & \textbf{.027} & \textbf{.562} & 3.577 & .632 & \textbf{.373} \\
共享校准RF（非LLM诊断） & .533 & \textbf{.491} & \textbf{.486} & .229 & .706 & .505 & .033 & .566 & \textbf{3.454} & \textbf{.649} & .383 \\
\bottomrule
\end{tabular}
}
\end{table}


\chapter{Qwen3-8B的DocFinQA结果}
\label{app:qwen}

Qwen3-8B在DocFinQA上返回解释性文字而非短答案，精确匹配接近地板。表\ref{tab:qwen-doc}报告该基座的全部条件；它不用于义务检索分支的主比较，只作为基座能力边界。该结果说明，短答案抽取任务对阅读器的格式遵循能力有额外要求；当阅读器返回解释性文字时，精确匹配和词级F1会同时下降，但答案覆盖率仍可能较高。

\begin{table}[htbp]
\centering
\zihao{5}
\setlength{\tabcolsep}{4pt}
\caption{Qwen3-8B的DocFinQA结果（300题固定样本）。}
\label{tab:qwen-doc}
\begin{tabular}{lrrrr}
\toprule
条件 & 准确率 & 宏F1 & 平均段落 & 覆盖率 \\
\midrule
固定4段 & .100 & .048 & 4.00 & .760 \\
固定8段 & .110 & .040 & 8.00 & .833 \\
固定16段 & .100 & .046 & 16.00 & .853 \\
规划器B=8 & .093 & .034 & 3.67 & .883 \\
规划器+回退 & .093 & .034 & 5.85 & .883 \\
\bottomrule
\end{tabular}
\end{table}

\chapter{命题与完整证明}
\label{app:proofs}

本章给出正文使用的四个命题及完整证明。证明只使用有限结果空间、有限动作集和线性期望，不依赖具体数据或模型。

\section{裁决门的贝叶斯最优性}

\begin{proposition}
给定信念$b\in\Delta(\mathcal H)$、动作集合$\mathcal A=\{\text{answer},\text{adjudicate},\text{abstain}\}$和损失函数$L:\mathcal A\times\mathcal H\to\mathbb R$，选择
\[
\pi^\star=\arg\min_{\pi\in\mathcal A}\mathbb E_{y\sim b}[L(\pi,y)]
\]
是贝叶斯最优动作。
\end{proposition}

\begin{proof}
对任意动作$\pi$，其在信念$b$下的期望损失为
\[
R(\pi,b)=\sum_{y\in\mathcal H}b(y)L(\pi,y).
\]
这是$b$的线性函数，也是有限集合$\mathcal A$上的实值函数。由$\pi^\star$的定义，
\[
R(\pi^\star,b)\le R(\pi,b),\qquad \forall \pi\in\mathcal A.
\]
因此不存在在$b$下期望损失更低的动作。若存在另一个动作$\pi'$满足$R(\pi',b)=R(\pi^\star,b)$，则$\pi'$与$\pi^\star$并列为贝叶斯最优；本文按固定顺序打破平局，不改变最优性。回答动作内部还要在三个结果桶中选择一个标签，但该选择同样是在有限集合上取最小期望损失，因而可以与动作选择合并为同一个贝叶斯决策。
\end{proof}

\section{路由效用的分解}

\begin{proposition}
设确定查询集合为$\mathcal D$，前瞻性查询集合为$\mathcal V$，路由器预测为$\hat{\mathcal D}$和$\hat{\mathcal V}$。混合效用可分解为
\[
U=\frac{1}{n}\sum_{q\in\mathcal D\cap\hat{\mathcal D}}u_D(q)
+\frac{1}{n}\sum_{q\in\mathcal V\cap\hat{\mathcal V}}u_V(q),
\]
其中$u_D(q)$为确定查询的F1，$u_V(q)=1-L(q)/10$为前瞻性查询的效用。
\end{proposition}

\begin{proof}
设$n=|\mathcal D|+|\mathcal V|$。对每条查询定义指示变量
\[
\mathbb I_D(q)=\mathbf 1[q\in\mathcal D],\qquad
\mathbb I_V(q)=\mathbf 1[q\in\mathcal V].
\]
路由器正确时获得对应分支效用，错误时获得0，因此
\[
U=\frac{1}{n}\sum_{q\in\mathcal D}\mathbb I[q\in\hat{\mathcal D}]u_D(q)
+\frac{1}{n}\sum_{q\in\mathcal V}\mathbb I[q\in\hat{\mathcal V}]u_V(q).
\]
把指示变量写成集合交集，即得
\[
U=\frac{1}{n}\sum_{q\in\mathcal D\cap\hat{\mathcal D}}u_D(q)
+\frac{1}{n}\sum_{q\in\mathcal V\cap\hat{\mathcal V}}u_V(q).
\]
oracle路由器对每条查询都正确，故
\[
U_{\text{oracle}}=\frac{1}{n}\sum_{q\in\mathcal D}u_D(q)+\frac{1}{n}\sum_{q\in\mathcal V}u_V(q).
\]
两者之差为错路由项，等于被错路由查询在本分支上的效用之和除以$n$。这一分解说明混合效用可以拆成“路由正确时的分支效用”和“错路由损失”两部分，从而把路由错误与分支错误分开诊断。
\end{proof}

\section{软锚点更新的归一化与单调性}

\begin{proposition}
设$\mathcal H$为有限结果空间，软锚点的隐含桶为$j$，似然为
\[
L_e(k)=
\begin{cases}
p_e, & k=j,\\
(1-p_e)/(|\mathcal H|-1), & k\ne j,
\end{cases}
\]
其中$p_e\in(1/|\mathcal H|,1)$。对任意先验$b(k)>0$和可靠性$\rho\in[0,1]$，更新
\[
b'(k)=\frac{b(k)L_e(k)^\rho}{\sum_{h\in\mathcal H}b(h)L_e(h)^\rho}
\]
是概率分布；当$\rho=0$时$b'=b$，当$\rho=1$时$b'$等于先验与完整似然的归一化乘积；$b'(j)/b'(k)$随$\rho$单调增大。
\end{proposition}

\begin{proof}
分母为正，因为$b(k)>0$且$L_e(k)>0$，所以$b'$非负且和为1，是概率分布。$\rho=0$时$L_e(k)^0=1$，故$b'=b$。对$k\ne j$，
\[
\frac{b'(j)}{b'(k)}
=\frac{b(j)}{b(k)}
\left(\frac{p_e}{(1-p_e)/(|\mathcal H|-1)}\right)^\rho.
\]
括号内的比值大于1（因为$p_e>1/|\mathcal H|$），所以该比值随$\rho$单调增大。可靠性越高，更新越偏向隐含桶；可靠性为0时锚点不改变信念。硬锚点可以看作$\rho\to\infty$并将非可行桶的似然置零的极限，本文用点质量加小底实现，以避免后续软锚点无法恢复。
\end{proof}

\section{可信集的最小基数}

\begin{proposition}
设$b$为$\mathcal H$上的概率分布，$R_\alpha$为满足$\sum_{k\in R}b(k)\ge 1-\alpha$的最小基数子集。把$b$按降序排列为$b_{(1)}\ge\cdots\ge b_{(|\mathcal H|)}$，则$R_\alpha$由前$m$个最大元素组成，其中$m$是使部分和达到$1-\alpha$的最小整数。
\end{proposition}

\begin{proof}
对任意大小为$m$的子集$R$，其概率质量不超过前$m$个最大元素之和。若前$m-1$个元素之和小于$1-\alpha$，则任何大小为$m-1$的子集都无法达到覆盖率，故最小基数至少为$m$。前$m$个元素之和达到$1-\alpha$，因此它们构成一个可行集，且基数等于下界$m$，所以是最小基数可信集。若存在并列的阈值元素，任选其中一个不改变覆盖率；本文按桶的固定顺序打破平局。
\end{proof}

\chapter{LLM提示词}
\label{app:prompts}

本章给出正文三个LLM调用的完整提示词。提示词只使用截止时点之前的输入，不包含已实现结果；JSON字段与正文的结构化输出一致。

\section{路由器提示词}

\begin{lstlisting}[basicstyle=\ttfamily\scriptsize,breaklines=true,columns=fullflexible,frame=single,numbers=none]
You decide whether a financial input is determinate or vague,
then route it to one of two specialists.

DETERMINATE (RAG): the input asks for a fact, number, passage,
estimate, or projection that a filing states or lets you compute
at the cutoff. A question about a future value is determinate when
the filing itself states that value or projection.

VAGUE (ANCHOR): the input asserts or asks about a future outcome
that is not realized at the cutoff. Examples include whether
reported revenue will fall below, within, or above a guidance
range, whether a product will launch, or whether a target will be
met. A management guidance statement is a forward-looking claim,
not a realized fact, even when it contains a number.

Return JSON only:
{"vague":true|false,"route":"rag|anchor",
 "confidence":0.0,"reason":"..."}

Examples:
Input: What was Entergy's net revenue change during 2015?
 -> {"vague":false,"route":"rag"}
Input: Q1 2024 revenue guidance is $12.2-13.2 billion.
 -> {"vague":true,"route":"anchor"}
Input: What did the 10-K report for segment operating profit?
 -> {"vague":false,"route":"rag"}
Input: Will Q3 revenue land within management's guidance range?
 -> {"vague":true,"route":"anchor"}
Input: Management expects second-half revenue of $1.5-1.6 billion.
 -> {"vague":true,"route":"anchor"}

Input: {query}
\end{lstlisting}

\section{锚点抽取提示词}

\begin{lstlisting}[basicstyle=\ttfamily\scriptsize,breaklines=true,columns=fullflexible,frame=single,numbers=none]
You are an AnchorEvolve adjudicator for a vague forward-looking
financial claim. Use ONLY information available at the cutoff.
Do not use the realized outcome.

The claim asks whether reported quarterly revenue will land below,
within, or above the guidance range. The guidance range defines the
three buckets; it is NOT by itself evidence that the outcome will be
within. Use the current-quarter value relative to the guidance
midpoint, the guidance width, and the observed base rates:
- mid_to_current_ratio > 1 means the midpoint requires sequential
  growth; the larger it is, the more likely below.
- mid_to_current_ratio < 1 means the midpoint implies a sequential
  decline; the smaller it is, the more likely above.
- A narrow guidance width makes within more likely; a very wide
  range makes outside buckets more likely.
- Base rates are priors, not destiny; update them with these features.

Extract explicit anchors and return JSON only:
{
  "anchors": [
    {"level":"L2",
     "type":"current_quarter_growth|guidance_width|sector_history|company_history|other",
     "relation":"bound|precedence|lead_time|duration|conservation|rate|unit_conversion|scope",
     "value":"...", "scope":"company|sector|market",
     "reliability":0.0, "hard":false,
     "implied_bucket":"below|within|above|null",
     "rationale":"..."}
  ],
  "bucket_scores":{"below":0.0,"within":0.0,"above":0.0},
  "rationale":"..."
}

Rules: a hard anchor is an explicit constraint that can rule out a
bucket (e.g. a reliable current-quarter value far above the guidance
range); most management guidance is a soft signal. reliability is in
[0,1]. bucket_scores must sum to 1.

Input:
{claim, company, ticker, sector, cutoff, target_year,
 target_quarter, guidance_low_musd, guidance_high_musd,
 guidance_mid_musd, guidance_width_pct, current_q_musd,
 current_q_reliability, observed_bucket_prior, sector_bucket_prior,
 evidence}
\end{lstlisting}

\section{阅读器提示词}

\begin{lstlisting}[basicstyle=\ttfamily\scriptsize,breaklines=true,columns=fullflexible,frame=single,numbers=none]
System:
You are a financial analyst answering a question about a long SEC
filing. Use only the excerpts below. The question may require
arithmetic; do the arithmetic carefully. If the excerpts do not
contain enough information to answer, return "insufficient" and set
sufficient to false.

Return only a JSON object with these keys:
- answer: the final answer as a short string (number or phrase)
- confidence: a number between 0 and 1
- sufficient: true if the excerpts are enough to answer, otherwise false
- missing: a short description of the missing information
- evidence_ids: the paragraph ids you used
- rationale: at most 40 words, no markdown

User:
Use only the excerpts below. Do not use outside knowledge.
This is acquisition round {round}; the excerpts include additional
paragraphs retrieved after you reported missing information.
Question: {question}
Document excerpts:
{evidence}
Return the JSON object now.
\end{lstlisting}

\section{提示词设计说明}

三个提示词共同遵循三条约束。第一，时点约束：路由器只接收查询与截止时点元数据，锚点提取器只接收截止时点之前的证据，阅读器只接收已获取的段落块。第二，结构化输出：三个调用都返回固定字段的JSON，解析失败在评价中计为错误，不进行静默修复。第三，可审计性：锚点提示词要求逐条给出类型、关系、作用域、可靠性、硬软标记、隐含桶和理由；阅读器要求给出证据编号和缺失信息，使答案可以追溯到具体段落或表格块。

\chapter{核心代码}
\label{app:code}

本章给出与正文算法直接对应的核心代码片段，完整实现和运行脚本随论文发布。代码只依赖公开的Python包和本地BM25索引，不包含API密钥。

\section{义务项与覆盖}

\begin{lstlisting}[language=Python]
import re

OBLIGATION_RE = re.compile(
    r"\$?\d[\d,]*(?:\.\d+)?%?|"
    r"\b(?:January|February|March|April|May|June|July|"
    r"August|September|October|November|December)\b",
    flags=re.I,
)

def obligation_terms(text):
    terms, seen = [], set()
    for match in OBLIGATION_RE.findall(str(text or "")):
        term = match.strip().lower().replace("$", "").replace(",", "").replace("%", "")
        if term in {"", "0", "1", "2", "3", "4", "5", "6", "7", "8", "9"}:
            continue
        if term not in seen:
            seen.add(term)
            terms.append(term)
    return terms

def uncovered_terms(terms, chunks):
    text = " ".join(str(getattr(c, "text", "")) for c in chunks)
    lowered = text.lower().replace(",", "")
    covered = set()
    for term in terms:
        if term.isdigit() and len(term) == 4 and term.startswith(("19", "20")):
            if re.search(r"(?<!\d)" + re.escape(term) + r"(?!\d)", lowered):
                covered.add(term)
        elif re.search(r"(?<![0-9.])" + re.escape(term) + r"(?![0-9])", lowered):
            covered.add(term)
    return [t for t in terms if t not in covered]
\end{lstlisting}

\section{锚点信念更新与裁决门}

\begin{lstlisting}[language=Python]
import numpy as np

BUCKETS = ["below", "within", "above"]
HARD_FLOOR = 1e-3
SOFT_P_E = 0.7
LOSS = {"over": 12.0, "under": 8.0, "wrong": 10.0,
        "miss": 8.0, "width": 2.0, "abstain": 5.0}

def normalize(v):
    v = np.maximum(np.asarray(v, dtype=float), 1e-12)
    return v / v.sum()

def belief_from_anchors(prior, parsed):
    b, anchors = prior.copy(), []
    if not parsed:
        return b, anchors
    for anchor in parsed.get("anchors", []) or []:
        implied = str(anchor.get("implied_bucket", "")).lower()
        if implied not in BUCKETS:
            continue
        rel = max(0.0, min(1.0, float(anchor.get("reliability", 0.5) or 0.5)))
        target = BUCKETS.index(implied)
        hard = bool(anchor.get("hard", False)) and rel >= 0.65
        anchors.append({"type": anchor.get("type"),
                        "relation": anchor.get("relation"),
                        "implied_bucket": implied,
                        "reliability": rel, "hard": hard,
                        "scope": anchor.get("scope")})
        if hard:
            b = np.array([1.0 if i == target else HARD_FLOOR for i in range(3)])
        else:
            like = np.full(3, (1.0 - SOFT_P_E) / 2.0)
            like[target] = SOFT_P_E
            b = b * np.power(like, rel)
        b = normalize(b)
    scores = parsed.get("bucket_scores")
    if isinstance(scores, dict):
        marginal = normalize(np.array([float(scores.get(x, 0.0)) for x in BUCKETS]))
        b = normalize(0.5 * b + 0.5 * marginal)
    return b, anchors

def credible_set(b, alpha=0.9):
    order = np.argsort(-b)
    out, total = [], 0.0
    for idx in order:
        out.append(int(idx))
        total += float(b[int(idx)])
        if total >= alpha:
            break
    return sorted(out)

def gate(b, alpha=0.9):
    exp_ans = [sum(float(b[y]) * (0.0 if a == y else
               LOSS["over"] if (a, y) == (2, 0) else
               LOSS["under"] if (a, y) == (0, 2) else
               LOSS["wrong"]) for y in range(3)) for a in range(3)]
    a = int(np.argmin(exp_ans))
    R = credible_set(b, alpha)
    exp_set = sum(float(b[y]) * ((0.0 if y in R else LOSS["miss"]) +
                 LOSS["width"] * max(0, len(R) - 1)) for y in range(3))
    if exp_ans[a] <= exp_set and exp_ans[a] <= LOSS["abstain"]:
        return {"mode": "answer", "a": a, "R": R}
    if exp_set <= LOSS["abstain"]:
        return {"mode": "adjudicate", "a": a, "R": R}
    return {"mode": "abstain", "a": a, "R": R}
\end{lstlisting}

\section{校准随机森林}

\begin{lstlisting}[language=Python]
import numpy as np
from sklearn.ensemble import RandomForestClassifier

for fold in range(FOLDS):
    train_rows = [row for row in rows
                  if fold_of(company_key(row[2]), seed) != fold
                  and row[1] < boundary]
    fold_test = [row for row in test_rows
                 if fold_of(company_key(row[2]), seed) == fold]
    X_train, C_train, y_train, meta_train = build(train_rows, by_company)
    X_test, C_test, y_test, meta_test = build(fold_test, by_company)
    train_features = np.hstack([X_train, C_train])
    test_features = np.hstack([X_test, C_test])
    model = make_rf(seed, class_weight, min_samples_leaf)
    model.fit(train_features, list(map(str, y_train)))
    probabilities = align(model.predict_proba(test_features), model.classes_)
    groups = np.asarray([company_key(record) for record in meta_train])
    oof, truths = [], []
    labels = np.asarray(list(map(str, y_train)))
    for fit_index, cal_index in GroupKFold(n_splits=min(3, len(set(groups)))).split(
        train_features, labels, groups
    ):
        inner = make_rf(seed, class_weight, min_samples_leaf)
        inner.fit(train_features[fit_index], labels[fit_index].tolist())
        oof.append(align(inner.predict_proba(train_features[cal_index]),
                         inner.classes_))
        truths.extend(labels[cal_index].tolist())
    temperature = fit_temperature(np.vstack(oof), truths)
\end{lstlisting}

\chapter{实现结构与复现边界}
\label{app:implementation}

\section{实现结构}

实现分为五个模块。路由模块读取原始查询和截止时点元数据，返回路由标签、置信度和简短理由。锚点模块读取断言、指引区间、当前季度值、基率和证据项，返回L2锚点及其可靠性、硬软标记和隐含桶。义务检索分支模块建立时点段落与表格索引，按义务覆盖迭代取证，并调用阅读器返回闭合答案。校准模块在开发集与验证集上训练随机森林并做温度缩放。汇总模块把逐条预测、路由决策和指标写入统一结果文件。各模块通过统一的数据结构连接，LLM调用按请求负载的哈希缓存，因此相同输入可以复现相同输出。

\section{复现边界}

公开结果包含逐条预测、路由决策和汇总指标；LLM响应缓存不随论文提交。随机森林的折种子、校准方式和融合权重选择在正文中给出。所有API密钥从环境变量读取，不写入结果文件或论文源文件。具名基线为同基座适配，代码、提示词和逐条输出随论文发布；原始系统在其原始数据上的未修改端到端结果不在本文的复现范围内。


% ==================== 致谢（终稿恢复） ====================
% 按学院通知：预答辩与查重阶段不排入致谢，终稿阶段取消注释并补写致谢内容。
% \clearpage
% \sufehead{致\quad 谢}

% ==================== 个人简历（终稿恢复） ====================
% 按学院通知：预答辩与查重阶段不排入个人简历及在学期间发表的研究成果，终稿阶段取消注释。
% \clearpage
% \sufehead{个人简历及在学期间发表的研究成果}
% {\zihao{5}\songti
% \setlength{\baselineskip}{16pt}
% \noindent 姓名：高深深\\
% 出生年月：\underline{\hspace{2.4cm}}年\underline{\hspace{1.2cm}}月\underline{\hspace{1.2cm}}日\\
% 本科毕业院校：\underline{\hspace{6cm}}\\
% 本科专业：\underline{\hspace{6cm}}\\
% 获学士学位时间：\underline{\hspace{2.4cm}}年\underline{\hspace{1.2cm}}月\\[0.4em]
% 硕士研究生：上海财经大学金融学院，金融工程与量化投资专业，2025年9月入学，导师罗丹。\\[0.6em]
% \textbf{在学期间发表的研究成果：}无。
% }

% ==================== 封底 ====================
\clearpage
\thispagestyle{empty}
\vspace*{\fill}
\begin{center}
{\zihao{3}\heiti 上海财经大学}\\[1.2em]
{\zihao{4}\songti 上海市国定路 777 号}
\end{center}
\vspace*{\fill}


\end{document}
